% concatenated from arXiv_v4.tex, math_commands.tex

\documentclass[10pt]{article} % For LaTeX2e
%\usepackage{tmlr}
% If accepted, instead use the following line for the camera-ready submission:
%\usepackage[preprint]{tmlr}
% To de-anonymize and remove mentions to TMLR (for example for posting to preprint servers), instead use the following:
%\usepackage[preprint]{tmlr}

% Optional math commands from https://github.com/goodfeli/dlbook_notation.
%%%%% NEW MATH DEFINITIONS %%%%%

\usepackage{amsmath,amsfonts,bm}

% Mark sections of captions for referring to divisions of figures
\newcommand{\figleft}{{\em (Left)}}
\newcommand{\figcenter}{{\em (Center)}}
\newcommand{\figright}{{\em (Right)}}
\newcommand{\figtop}{{\em (Top)}}
\newcommand{\figbottom}{{\em (Bottom)}}
\newcommand{\captiona}{{\em (a)}}
\newcommand{\captionb}{{\em (b)}}
\newcommand{\captionc}{{\em (c)}}
\newcommand{\captiond}{{\em (d)}}

% Highlight a newly defined term
\newcommand{\newterm}[1]{{\bf #1}}


% Figure reference, lower-case.
\def\figref#1{figure~\ref{#1}}
% Figure reference, capital. For start of sentence
\def\Figref#1{Figure~\ref{#1}}
\def\twofigref#1#2{figures \ref{#1} and \ref{#2}}
\def\quadfigref#1#2#3#4{figures \ref{#1}, \ref{#2}, \ref{#3} and \ref{#4}}
% Section reference, lower-case.
\def\secref#1{section~\ref{#1}}
% Section reference, capital.
\def\Secref#1{Section~\ref{#1}}
% Reference to two sections.
\def\twosecrefs#1#2{sections \ref{#1} and \ref{#2}}
% Reference to three sections.
\def\secrefs#1#2#3{sections \ref{#1}, \ref{#2} and \ref{#3}}
% Reference to an equation, lower-case.
\def\eqref#1{equation~\ref{#1}}
% Reference to an equation, upper case
\def\Eqref#1{Equation~\ref{#1}}
% A raw reference to an equation---avoid using if possible
\def\plaineqref#1{\ref{#1}}
% Reference to a chapter, lower-case.
\def\chapref#1{chapter~\ref{#1}}
% Reference to an equation, upper case.
\def\Chapref#1{Chapter~\ref{#1}}
% Reference to a range of chapters
\def\rangechapref#1#2{chapters\ref{#1}--\ref{#2}}
% Reference to an algorithm, lower-case.
\def\algref#1{algorithm~\ref{#1}}
% Reference to an algorithm, upper case.
\def\Algref#1{Algorithm~\ref{#1}}
\def\twoalgref#1#2{algorithms \ref{#1} and \ref{#2}}
\def\Twoalgref#1#2{Algorithms \ref{#1} and \ref{#2}}
% Reference to a part, lower case
\def\partref#1{part~\ref{#1}}
% Reference to a part, upper case
\def\Partref#1{Part~\ref{#1}}
\def\twopartref#1#2{parts \ref{#1} and \ref{#2}}

\def\ceil#1{\lceil #1 \rceil}
\def\floor#1{\lfloor #1 \rfloor}
\def\1{\bm{1}}
\newcommand{\train}{\mathcal{D}}
\newcommand{\valid}{\mathcal{D_{\mathrm{valid}}}}
\newcommand{\test}{\mathcal{D_{\mathrm{test}}}}

\def\eps{{\epsilon}}


% Random variables
\def\reta{{\textnormal{$\eta$}}}
\def\ra{{\textnormal{a}}}
\def\rb{{\textnormal{b}}}
\def\rc{{\textnormal{c}}}
\def\rd{{\textnormal{d}}}
\def\re{{\textnormal{e}}}
\def\rf{{\textnormal{f}}}
\def\rg{{\textnormal{g}}}
\def\rh{{\textnormal{h}}}
\def\ri{{\textnormal{i}}}
\def\rj{{\textnormal{j}}}
\def\rk{{\textnormal{k}}}
\def\rl{{\textnormal{l}}}
% rm is already a command, just don't name any random variables m
\def\rn{{\textnormal{n}}}
\def\ro{{\textnormal{o}}}
\def\rp{{\textnormal{p}}}
\def\rq{{\textnormal{q}}}
\def\rr{{\textnormal{r}}}
\def\rs{{\textnormal{s}}}
\def\rt{{\textnormal{t}}}
\def\ru{{\textnormal{u}}}
\def\rv{{\textnormal{v}}}
\def\rw{{\textnormal{w}}}
\def\rx{{\textnormal{x}}}
\def\ry{{\textnormal{y}}}
\def\rz{{\textnormal{z}}}

% Random vectors
\def\rvepsilon{{\mathbf{\epsilon}}}
\def\rvtheta{{\mathbf{\theta}}}
\def\rva{{\mathbf{a}}}
\def\rvb{{\mathbf{b}}}
\def\rvc{{\mathbf{c}}}
\def\rvd{{\mathbf{d}}}
\def\rve{{\mathbf{e}}}
\def\rvf{{\mathbf{f}}}
\def\rvg{{\mathbf{g}}}
\def\rvh{{\mathbf{h}}}
\def\rvu{{\mathbf{i}}}
\def\rvj{{\mathbf{j}}}
\def\rvk{{\mathbf{k}}}
\def\rvl{{\mathbf{l}}}
\def\rvm{{\mathbf{m}}}
\def\rvn{{\mathbf{n}}}
\def\rvo{{\mathbf{o}}}
\def\rvp{{\mathbf{p}}}
\def\rvq{{\mathbf{q}}}
\def\rvr{{\mathbf{r}}}
\def\rvs{{\mathbf{s}}}
\def\rvt{{\mathbf{t}}}
\def\rvu{{\mathbf{u}}}
\def\rvv{{\mathbf{v}}}
\def\rvw{{\mathbf{w}}}
\def\rvx{{\mathbf{x}}}
\def\rvy{{\mathbf{y}}}
\def\rvz{{\mathbf{z}}}

% Elements of random vectors
\def\erva{{\textnormal{a}}}
\def\ervb{{\textnormal{b}}}
\def\ervc{{\textnormal{c}}}
\def\ervd{{\textnormal{d}}}
\def\erve{{\textnormal{e}}}
\def\ervf{{\textnormal{f}}}
\def\ervg{{\textnormal{g}}}
\def\ervh{{\textnormal{h}}}
\def\ervi{{\textnormal{i}}}
\def\ervj{{\textnormal{j}}}
\def\ervk{{\textnormal{k}}}
\def\ervl{{\textnormal{l}}}
\def\ervm{{\textnormal{m}}}
\def\ervn{{\textnormal{n}}}
\def\ervo{{\textnormal{o}}}
\def\ervp{{\textnormal{p}}}
\def\ervq{{\textnormal{q}}}
\def\ervr{{\textnormal{r}}}
\def\ervs{{\textnormal{s}}}
\def\ervt{{\textnormal{t}}}
\def\ervu{{\textnormal{u}}}
\def\ervv{{\textnormal{v}}}
\def\ervw{{\textnormal{w}}}
\def\ervx{{\textnormal{x}}}
\def\ervy{{\textnormal{y}}}
\def\ervz{{\textnormal{z}}}

% Random matrices
\def\rmA{{\mathbf{A}}}
\def\rmB{{\mathbf{B}}}
\def\rmC{{\mathbf{C}}}
\def\rmD{{\mathbf{D}}}
\def\rmE{{\mathbf{E}}}
\def\rmF{{\mathbf{F}}}
\def\rmG{{\mathbf{G}}}
\def\rmH{{\mathbf{H}}}
\def\rmI{{\mathbf{I}}}
\def\rmJ{{\mathbf{J}}}
\def\rmK{{\mathbf{K}}}
\def\rmL{{\mathbf{L}}}
\def\rmM{{\mathbf{M}}}
\def\rmN{{\mathbf{N}}}
\def\rmO{{\mathbf{O}}}
\def\rmP{{\mathbf{P}}}
\def\rmQ{{\mathbf{Q}}}
\def\rmR{{\mathbf{R}}}
\def\rmS{{\mathbf{S}}}
\def\rmT{{\mathbf{T}}}
\def\rmU{{\mathbf{U}}}
\def\rmV{{\mathbf{V}}}
\def\rmW{{\mathbf{W}}}
\def\rmX{{\mathbf{X}}}
\def\rmY{{\mathbf{Y}}}
\def\rmZ{{\mathbf{Z}}}

% Elements of random matrices
\def\ermA{{\textnormal{A}}}
\def\ermB{{\textnormal{B}}}
\def\ermC{{\textnormal{C}}}
\def\ermD{{\textnormal{D}}}
\def\ermE{{\textnormal{E}}}
\def\ermF{{\textnormal{F}}}
\def\ermG{{\textnormal{G}}}
\def\ermH{{\textnormal{H}}}
\def\ermI{{\textnormal{I}}}
\def\ermJ{{\textnormal{J}}}
\def\ermK{{\textnormal{K}}}
\def\ermL{{\textnormal{L}}}
\def\ermM{{\textnormal{M}}}
\def\ermN{{\textnormal{N}}}
\def\ermO{{\textnormal{O}}}
\def\ermP{{\textnormal{P}}}
\def\ermQ{{\textnormal{Q}}}
\def\ermR{{\textnormal{R}}}
\def\ermS{{\textnormal{S}}}
\def\ermT{{\textnormal{T}}}
\def\ermU{{\textnormal{U}}}
\def\ermV{{\textnormal{V}}}
\def\ermW{{\textnormal{W}}}
\def\ermX{{\textnormal{X}}}
\def\ermY{{\textnormal{Y}}}
\def\ermZ{{\textnormal{Z}}}

% Vectors
\def\vzero{{\bm{0}}}
\def\vone{{\bm{1}}}
\def\vmu{{\bm{\mu}}}
\def\vtheta{{\bm{\theta}}}
\def\va{{\bm{a}}}
\def\vb{{\bm{b}}}
\def\vc{{\bm{c}}}
\def\vd{{\bm{d}}}
\def\ve{{\bm{e}}}
\def\vf{{\bm{f}}}
\def\vg{{\bm{g}}}
\def\vh{{\bm{h}}}
\def\vi{{\bm{i}}}
\def\vj{{\bm{j}}}
\def\vk{{\bm{k}}}
\def\vl{{\bm{l}}}
\def\vm{{\bm{m}}}
\def\vn{{\bm{n}}}
\def\vo{{\bm{o}}}
\def\vp{{\bm{p}}}
\def\vq{{\bm{q}}}
\def\vr{{\bm{r}}}
\def\vs{{\bm{s}}}
\def\vt{{\bm{t}}}
\def\vu{{\bm{u}}}
\def\vv{{\bm{v}}}
\def\vw{{\bm{w}}}
\def\vx{{\bm{x}}}
\def\vy{{\bm{y}}}
\def\vz{{\bm{z}}}

% Elements of vectors
\def\evalpha{{\alpha}}
\def\evbeta{{\beta}}
\def\evepsilon{{\epsilon}}
\def\evlambda{{\lambda}}
\def\evomega{{\omega}}
\def\evmu{{\mu}}
\def\evpsi{{\psi}}
\def\evsigma{{\sigma}}
\def\evtheta{{\theta}}
\def\eva{{a}}
\def\evb{{b}}
\def\evc{{c}}
\def\evd{{d}}
\def\eve{{e}}
\def\evf{{f}}
\def\evg{{g}}
\def\evh{{h}}
\def\evi{{i}}
\def\evj{{j}}
\def\evk{{k}}
\def\evl{{l}}
\def\evm{{m}}
\def\evn{{n}}
\def\evo{{o}}
\def\evp{{p}}
\def\evq{{q}}
\def\evr{{r}}
\def\evs{{s}}
\def\evt{{t}}
\def\evu{{u}}
\def\evv{{v}}
\def\evw{{w}}
\def\evx{{x}}
\def\evy{{y}}
\def\evz{{z}}

% Matrix
\def\mA{{\bm{A}}}
\def\mB{{\bm{B}}}
\def\mC{{\bm{C}}}
\def\mD{{\bm{D}}}
\def\mE{{\bm{E}}}
\def\mF{{\bm{F}}}
\def\mG{{\bm{G}}}
\def\mH{{\bm{H}}}
\def\mI{{\bm{I}}}
\def\mJ{{\bm{J}}}
\def\mK{{\bm{K}}}
\def\mL{{\bm{L}}}
\def\mM{{\bm{M}}}
\def\mN{{\bm{N}}}
\def\mO{{\bm{O}}}
\def\mP{{\bm{P}}}
\def\mQ{{\bm{Q}}}
\def\mR{{\bm{R}}}
\def\mS{{\bm{S}}}
\def\mT{{\bm{T}}}
\def\mU{{\bm{U}}}
\def\mV{{\bm{V}}}
\def\mW{{\bm{W}}}
\def\mX{{\bm{X}}}
\def\mY{{\bm{Y}}}
\def\mZ{{\bm{Z}}}
\def\mBeta{{\bm{\beta}}}
\def\mPhi{{\bm{\Phi}}}
\def\mLambda{{\bm{\Lambda}}}
\def\mSigma{{\bm{\Sigma}}}

% Tensor
\DeclareMathAlphabet{\mathsfit}{\encodingdefault}{\sfdefault}{m}{sl}
\SetMathAlphabet{\mathsfit}{bold}{\encodingdefault}{\sfdefault}{bx}{n}
\newcommand{\tens}[1]{\bm{\mathsfit{#1}}}
\def\tA{{\tens{A}}}
\def\tB{{\tens{B}}}
\def\tC{{\tens{C}}}
\def\tD{{\tens{D}}}
\def\tE{{\tens{E}}}
\def\tF{{\tens{F}}}
\def\tG{{\tens{G}}}
\def\tH{{\tens{H}}}
\def\tI{{\tens{I}}}
\def\tJ{{\tens{J}}}
\def\tK{{\tens{K}}}
\def\tL{{\tens{L}}}
\def\tM{{\tens{M}}}
\def\tN{{\tens{N}}}
\def\tO{{\tens{O}}}
\def\tP{{\tens{P}}}
\def\tQ{{\tens{Q}}}
\def\tR{{\tens{R}}}
\def\tS{{\tens{S}}}
\def\tT{{\tens{T}}}
\def\tU{{\tens{U}}}
\def\tV{{\tens{V}}}
\def\tW{{\tens{W}}}
\def\tX{{\tens{X}}}
\def\tY{{\tens{Y}}}
\def\tZ{{\tens{Z}}}


% Graph
\def\gA{{\mathcal{A}}}
\def\gB{{\mathcal{B}}}
\def\gC{{\mathcal{C}}}
\def\gD{{\mathcal{D}}}
\def\gE{{\mathcal{E}}}
\def\gF{{\mathcal{F}}}
\def\gG{{\mathcal{G}}}
\def\gH{{\mathcal{H}}}
\def\gI{{\mathcal{I}}}
\def\gJ{{\mathcal{J}}}
\def\gK{{\mathcal{K}}}
\def\gL{{\mathcal{L}}}
\def\gM{{\mathcal{M}}}
\def\gN{{\mathcal{N}}}
\def\gO{{\mathcal{O}}}
\def\gP{{\mathcal{P}}}
\def\gQ{{\mathcal{Q}}}
\def\gR{{\mathcal{R}}}
\def\gS{{\mathcal{S}}}
\def\gT{{\mathcal{T}}}
\def\gU{{\mathcal{U}}}
\def\gV{{\mathcal{V}}}
\def\gW{{\mathcal{W}}}
\def\gX{{\mathcal{X}}}
\def\gY{{\mathcal{Y}}}
\def\gZ{{\mathcal{Z}}}

% Sets
\def\sA{{\mathbb{A}}}
\def\sB{{\mathbb{B}}}
\def\sC{{\mathbb{C}}}
\def\sD{{\mathbb{D}}}
% Don't use a set called E, because this would be the same as our symbol
% for expectation.
\def\sF{{\mathbb{F}}}
\def\sG{{\mathbb{G}}}
\def\sH{{\mathbb{H}}}
\def\sI{{\mathbb{I}}}
\def\sJ{{\mathbb{J}}}
\def\sK{{\mathbb{K}}}
\def\sL{{\mathbb{L}}}
\def\sM{{\mathbb{M}}}
\def\sN{{\mathbb{N}}}
\def\sO{{\mathbb{O}}}
\def\sP{{\mathbb{P}}}
\def\sQ{{\mathbb{Q}}}
\def\sR{{\mathbb{R}}}
\def\sS{{\mathbb{S}}}
\def\sT{{\mathbb{T}}}
\def\sU{{\mathbb{U}}}
\def\sV{{\mathbb{V}}}
\def\sW{{\mathbb{W}}}
\def\sX{{\mathbb{X}}}
\def\sY{{\mathbb{Y}}}
\def\sZ{{\mathbb{Z}}}

% Entries of a matrix
\def\emLambda{{\Lambda}}
\def\emA{{A}}
\def\emB{{B}}
\def\emC{{C}}
\def\emD{{D}}
\def\emE{{E}}
\def\emF{{F}}
\def\emG{{G}}
\def\emH{{H}}
\def\emI{{I}}
\def\emJ{{J}}
\def\emK{{K}}
\def\emL{{L}}
\def\emM{{M}}
\def\emN{{N}}
\def\emO{{O}}
\def\emP{{P}}
\def\emQ{{Q}}
\def\emR{{R}}
\def\emS{{S}}
\def\emT{{T}}
\def\emU{{U}}
\def\emV{{V}}
\def\emW{{W}}
\def\emX{{X}}
\def\emY{{Y}}
\def\emZ{{Z}}
\def\emSigma{{\Sigma}}

% entries of a tensor
% Same font as tensor, without \bm wrapper
\newcommand{\etens}[1]{\mathsfit{#1}}
\def\etLambda{{\etens{\Lambda}}}
\def\etA{{\etens{A}}}
\def\etB{{\etens{B}}}
\def\etC{{\etens{C}}}
\def\etD{{\etens{D}}}
\def\etE{{\etens{E}}}
\def\etF{{\etens{F}}}
\def\etG{{\etens{G}}}
\def\etH{{\etens{H}}}
\def\etI{{\etens{I}}}
\def\etJ{{\etens{J}}}
\def\etK{{\etens{K}}}
\def\etL{{\etens{L}}}
\def\etM{{\etens{M}}}
\def\etN{{\etens{N}}}
\def\etO{{\etens{O}}}
\def\etP{{\etens{P}}}
\def\etQ{{\etens{Q}}}
\def\etR{{\etens{R}}}
\def\etS{{\etens{S}}}
\def\etT{{\etens{T}}}
\def\etU{{\etens{U}}}
\def\etV{{\etens{V}}}
\def\etW{{\etens{W}}}
\def\etX{{\etens{X}}}
\def\etY{{\etens{Y}}}
\def\etZ{{\etens{Z}}}

% The true underlying data generating distribution
\newcommand{\pdata}{p_{\rm{data}}}
% The empirical distribution defined by the training set
\newcommand{\ptrain}{\hat{p}_{\rm{data}}}
\newcommand{\Ptrain}{\hat{P}_{\rm{data}}}
% The model distribution
\newcommand{\pmodel}{p_{\rm{model}}}
\newcommand{\Pmodel}{P_{\rm{model}}}
\newcommand{\ptildemodel}{\tilde{p}_{\rm{model}}}
% Stochastic autoencoder distributions
\newcommand{\pencode}{p_{\rm{encoder}}}
\newcommand{\pdecode}{p_{\rm{decoder}}}
\newcommand{\precons}{p_{\rm{reconstruct}}}

\newcommand{\laplace}{\mathrm{Laplace}} % Laplace distribution

\newcommand{\E}{\mathbb{E}}
\newcommand{\Ls}{\mathcal{L}}
\newcommand{\R}{\mathbb{R}}
\newcommand{\emp}{\tilde{p}}
\newcommand{\lr}{\alpha}
\newcommand{\reg}{\lambda}
\newcommand{\rect}{\mathrm{rectifier}}
\newcommand{\softmax}{\mathrm{softmax}}
\newcommand{\sigmoid}{\sigma}
\newcommand{\softplus}{\zeta}
\newcommand{\KL}{D_{\mathrm{KL}}}
\newcommand{\Var}{\mathrm{Var}}
\newcommand{\standarderror}{\mathrm{SE}}
\newcommand{\Cov}{\mathrm{Cov}}
% Wolfram Mathworld says $L^2$ is for function spaces and $\ell^2$ is for vectors
% But then they seem to use $L^2$ for vectors throughout the site, and so does
% wikipedia.
\newcommand{\normlzero}{L^0}
\newcommand{\normlone}{L^1}
\newcommand{\normltwo}{L^2}
\newcommand{\normlp}{L^p}
\newcommand{\normmax}{L^\infty}

\newcommand{\parents}{Pa} % See usage in notation.tex. Chosen to match Daphne's book.

\DeclareMathOperator*{\argmax}{arg\,max}
\DeclareMathOperator*{\argmin}{arg\,min}

\DeclareMathOperator{\sign}{sign}
\DeclareMathOperator{\Tr}{Tr}
\let\ab\allowbreak


\usepackage{hyperref}
\usepackage{url}
\usepackage{graphicx}
\usepackage{float}
\usepackage{placeins}

\ExplSyntaxOn
\bool_new:N \accepted % Declare a boolean variable
%\bool_set_true:N \accepted % Set to true; 
\bool_set_false:N \accepted
\ExplSyntaxOff

% Attempt to make hyperref and algorithmic work together better:
\newcommand{\theHalgorithm}{\arabic{algorithm}}


\ExplSyntaxOn
\bool_if:NTF \accepted
    {\usepackage[accepted]{tmlr}}
    {\usepackage[preprint]{tmlr}}
\ExplSyntaxOff



% For theorems and such
\usepackage{amsmath}
\usepackage{amssymb}
\usepackage{mathtools}
\usepackage{amsthm}
\usepackage{amsfonts}
\usepackage{mathrsfs}
\usepackage{subcaption}

% if you use cleveref..
\usepackage[capitalize,noabbrev]{cleveref}

%%%%%%%%%%%%%%%%%%%%%%%%%%%%%%%%
% THEOREMS
%%%%%%%%%%%%%%%%%%%%%%%%%%%%%%%%
\theoremstyle{plain}
\newtheorem{theorem}{Theorem}[section]
\newtheorem{proposition}[theorem]{Proposition}
\newtheorem{lemma}[theorem]{Lemma}
\newtheorem{corollary}[theorem]{Corollary}
\theoremstyle{definition}
\newtheorem{definition}[theorem]{Definition}
\newtheorem{assumption}[theorem]{Assumption}
\theoremstyle{remark}
\newtheorem{remark}[theorem]{Remark}
\newtheorem{example}{Example}

% Todonotes is useful during development; simply uncomment the next line
%    and comment out the line below the next line to turn off comments
%\usepackage[disable,textsize=tiny]{todonotes}
%\usepackage[textsize=tiny]{todonotes}

%%%% Our own packages
%\newcommand{\E}{{\mathbb E}}
\renewcommand{\P}{{\mathbb P}}
%\newcommand{\R}{{\mathbb R}}
\renewcommand{\L}{\mathscr{L}}
\newcommand{\W}{\mathcal{W}}
\newcommand{\tr}{{\rm tr}}
\newcommand{\norm}[1]{\left\|#1\right\|}
\newcommand{\opnorm}[2][H]{\left\|#2\right\|_{\mathcal{L}(H, #1)}}
\newcommand{\C}{\mathcal C}
\newcommand{\N}{\mathcal N}
\newcommand{\Natural}{\mathbb N}
\newcommand{\cl}{\text{cl}}
\newcommand{\low}[1]{{\lfloor #1 \rfloor}}
\newcommand{\var}{\operatorname{Var}} 
\newcommand{\dd}{{\mathrm d}}

\newcommand{\fabian}[1]{{#1}}%{{\color{blue} #1}}
\newcommand{\note}[1]{#1}%{{\color{red} #1}}
\definecolor{darkgreen}{rgb}{0,0.5,0}
\newcommand{\lam}[1]{{ #1}}
\newcommand{\duclam}[1]{#1}%{{\color{darkgreen} #1}}
\newcommand{\tapio}[1]{#1}%{{\textcolor{red}{#1}\index{#1}}}
%%%%




\title{An Unconditional Representation of the Conditional Score in Infinite-Dimensional Linear Inverse Problems}
% alternative: Scalable task-dependent score diffusion posterior sampling in infinite-dimensional
% Authors must not appear in the submitted version. They should be hidden
% as long as the tmlr package is used without the [accepted] or [preprint] options.
% Non-anonymous submissions will be rejected without review.

\author{\name Fabian Schneider \email fabian.schneider@lut.fi \\
      \addr School of Engineering Science\\
      Lappeenranta-Lahti University of Technology \\%\\Finland
      Vienna University of Technology (TU Wien)
      \AND
      \name Duc-Lam Duong \email duc-lam.duong@lut.fi \\
      \addr School of Engineering Science\\
      Lappeenranta-Lahti University of Technology %\\Finland
      \AND
      \name Matti Lassas \email matti.lassas@helsinki.fi\\
      \addr Department of Mathematics and Statistics\\
      University of Helsinki %\\Finland
      \AND 
      \name Maarten V. de Hoop \email mdehoop@rice.edu \\
      \addr Department of Computational and Applied Mathematics\\
      Rice University %\\United States 
      \AND
      \name Tapio Helin \email tapio.helin@lut.fi \\
      \addr School of Engineering Science\\
      Lappeenranta-Lahti University of Technology %\\Finland
      }

% The \author macro works with any number of authors. Use \AND 
% to separate the names and addresses of multiple authors.

\newcommand{\fix}{\marginpar{FIX}}
\newcommand{\new}{\marginpar{NEW}}



\begin{document}
\def\month{11}  % Insert correct month for camera-ready version
\def\year{2025} % Insert correct year for camera-ready version
\def\openreview{\url{https://openreview.net/forum?id=rO8erhXHPo}} % Insert correct link to OpenReview for camera-ready version


\maketitle

\begin{abstract}
Score-based diffusion models (SDMs) have emerged as a powerful tool for sampling from the posterior distribution in Bayesian inverse problems. However, existing methods often require multiple evaluations of the forward mapping to generate a single sample, resulting in significant computational costs for large-scale inverse problems. To address this, we propose an unconditional representation of the conditional score function (UCoS) tailored to linear inverse problems, which avoids forward model evaluations during sampling by shifting computational effort to an offline training phase. In this phase, a \emph{task-dependent} score function is learned based on the linear forward operator. Crucially, we show that the conditional score can be derived \emph{exactly} from a trained (unconditional) score using affine transformations, eliminating the need for conditional score approximations. Our approach is formulated in infinite-dimensional function spaces, making it inherently discretization-invariant. We support this formulation with a rigorous convergence analysis that justifies UCoS beyond any specific discretization. Finally we validate UCoS through high-dimensional computed tomography (CT) and image deblurring experiments, demonstrating both scalability and accuracy.
\end{abstract}

\section{Introduction}
%\duclam{{23/5/25:} Lam editing and testing the TMLR \LaTeX\ format}\\

Inverse problems seek to determine unknown quantities through indirect and noisy measurements, typically leading to ill-posed scenarios.
The Bayesian approach to inverse problems frames the task as a quest for information. Blending statistical prior information of the unknown with a likelihood model for the measurement data gives rise to a posterior distribution, which fully characterizes the unknown conditioned on noisy data \cite{kaipio2006statistical, stuart2010inverse}. In severely ill-posed problems, the quality of inference is strongly dependent on the expressivity of the prior. Traditional hand-crafted priors, such as the total-variation prior, tend not to be expressive enough to characterize complicated structures \cite{sun2023provable}.
Generative models offer a flexible and computationally feasible approach to prior modeling as they offer the possibility of generating new samples after training on a data set characterizing the prior.

This work investigates sampling from the posterior distribution of linear inverse problems using score-based diffusion models (SDMs) \cite{song2021scorebasedDiffusion}, which have recently received wide attention in the literature (in the context of inverse problems, see e.g. \citet{batzolis2021conditional, lim2023score, hagemann2023multilevel, graikos2023diffusionplugandplay, feng2023scorebased, sun2023provable, pidstrigach2023infinitedimensional, Dey_2024, holzschuh2023solving, dou2023diffusion, barbano2023steerable, cardoso2023monte, feng2023efficient, song2024solving, meng2022diffusion, kveton2024online, wu2024practical, wu2024principled, baldassari2024taming, yao2025guideddiffusionsamplingfunction, chen2025solvinginverseproblemsdiffusionbased}). 
An SDM consists of two main components: a forward diffusion process and a reverse generative process. In the forward diffusion, the model gradually transforms the target distribution into a simpler, tractable distribution with a specific stochastic differential equation (SDE) driven by Gaussian noise. The generative process simulates the time-reversal of the diffusion process with a backwards SDE. This denoising phase relies on the drift, which is computed using the logarithmic gradients (scores) of the diffused data densities. These scores are typically estimated using a neural network, allowing efficient simulation of the backwards SDE and, consequently, sample generation from the target distribution. SDMs have demonstrated significant success across a variety of domains, including inverse problems such as medical imaging \cite{Dey_2024, levac2023mri, barbano2023steerable, chung2023solving, song2022solving}.

For inverse problems, posterior sampling with SDMs involves the challenging task of estimating the score function conditioned on the measurements. A key difficulty is balancing computational efficiency with scalability while ensuring accurate posterior sampling. Motivated by the limitations of existing approaches, we propose a novel method for conditional sampling, based on an \emph{unconditional task-dependent representation of the conditional score (UCoS)} that overcomes these challenges by leveraging an explicit likelihood model, which is often available and predefined in inverse problems. UCoS is particularly well-suited for large-scale imaging tasks that require efficient posterior sampling for varying sets of measurement data.


The existing posterior sampling methods can be broadly categorized into three approaches. The \emph{first} approach modifies the unconditional reverse diffusion process to guide sample trajectories towards the posterior, either by adding a correction term to the score \cite{jalal2021robustannealedlangevian, chung2023diffusion, song2023pseudoinverseguided, Adam2022Posteriorsamples, CHUNG2022102479, levac2023mri} or incorporating a data consistency optimization step \cite{song2022solving, Dey_2024, Chung2022ComeCloser}.
The \emph{second} approach employs gradient-based Monte Carlo sampling techniques, replacing the prior score with a learned score \cite{cardoso2023monte, sun2023provable, feng2023scorebased}, with recent extensions to infinite dimensions \cite{baldassari2024taming}. \duclam{For a recent review of methods largely following these two approaches (in finite dimensions), we refer to \citet{daras2024survey}.} A key advantage of these approaches is that they allow the use of a pre-trained unconditional prior score, accommodating different imaging tasks simultaneously without requiring task-specific training. This is particularly useful for handling extensive multipurpose training sets in applications such as image processing. However, such methods suffer from (1) high computational costs due to repeated forward evaluations, (2) poor scalability in high-dimensional or large-scale inverse problems, and (3) potential inconsistency in posterior samples due to the lack of rigorous convergence guarantees.

The \emph{third} approach directly trains a \emph{conditional} score function to perform conditional sampling, learning an amortized version of the conditional score that depends on observations \cite{baldassari2024conditional, batzolis2021conditional}. This method extends rigorously to infinite-dimensional diffusion models, making it attractive for large-scale inverse problems, as it eliminates the need to evaluate the forward map during sampling. However, it requires extensive training data on the joint distribution to reliably approximate the conditional score function, particularly as the problem dimensionality increases.

%Each of these three approaches has notable drawbacks. The first two \emph{task-independent} approaches suffer from (1) high computational costs due to repeated forward evaluations, (2) poor scalability in high-dimensional or large-scale inverse problems, and (3) potential inconsistency in posterior samples due to the lack of rigorous guarantees. While the third approach has theoretical guarantees, 

We address the prohibitive sampling cost for large-scale inverse problems by demonstrating that this computational overhead can be offloaded to the training phase, thereby\lam{improving the efficiency of posterior sampling}. This is achieved by designing a \emph{task-dependent} training phase that depends on the forward mapping but \emph{not} on the measurement data. More precisely, we introduce a diffusion-like random process whose distribution is explicitly dependent on the forward mapping, allowing the score (which we term the \emph{task-dependent score}) to be estimated using standard methods. We then demonstrate that the conditional score corresponding to the posterior distribution can be recovered from this task-dependent \emph{unconditional} score through simple affine transformations involving the measurement data. Furthermore, we modify the training procedure to enable evaluation of the conditional score without requiring application of the forward mapping during sampling.  In comparison to the conditional method, \lam{integrating the explicit likelihood model into the process reduces the dimensionality of the training target. As a result, UCoS offers the best of both worlds: \fabian{like the conditional approach,} it avoids forward mapping evaluations during posterior sampling while its score model remains as lightweight as standard unconditional ones. This leads to more favorable scaling of the online sampling cost as the size or dimensionality of the problem grows. }

Inverse problems are often about inferring quantities represented by functions such as solutions or parameters to PDEs \fabian{or integral equations. These models must be arbitrarily discretized to obtain a finite-dimensional computational model \cite{Mueller2012Linearandnonlinear}}. In Bayesian inversion, there has been a long-standing effort to develop methods that are discretization-independent \cite{lehtinen1989linear, lassas2004can, lassas2009}, aligned with the principle to ``avoid discretization until the last possible moment" \cite{stuart2010inverse}. This is also critical for the success of SDMs in Bayesian inversion as recent theoretical studies indicate that the performance guarantees do not always generalize well on increasing dimension \cite{chen2022sampling, de2022convergence, pidstrigach2023infinitedimensional}. Inspired by recent developments on defining the score-based diffusion framework in infinite-dimensional spaces \cite{pidstrigach2023infinitedimensional, baldassari2024conditional}, we also extend our method rigorously to a separable Hilbert space setting. In particular, in the spirit of \citet{pidstrigach2023infinitedimensional}, we perform a convergence analysis of UCoS, establishing a bound between the samples generated by UCoS and the true posterior target measure. Moreover, we conduct numerical experiments of inverse problems related to computerized tomography (CT) \fabian{and deblurring} to illustrate the practical applicability of UCoS.

\subsection{Related work}

The body of literature on SDMs is growing rapidly. Let us mention that \citet{song2021scorebasedDiffusion} developed a unified framework combining score-based \cite{hyvarinen2005estimation, song2019generative} and diffusion \cite{sohl2015deep, ho2020denoising} models to interpret SDMs as a time-reversal of certain SDEs. 
Our paper is inspired by the work on conditional SDMs in the context of inverse problems.
One line of work seeks to modify the unconditional prior score function to generate samples that approximately follow the posterior distribution. Examples of this include projection-based approaches \cite{song2022solving, Dey_2024}, which project the samples on a subspace that solves the inverse problem during the generating process;
methods based on the gradient of a log-likelihood \cite{jalal2021robustannealedlangevian, chung2023diffusion, levac2023mri} and plug and play approaches \cite{graikos2023diffusionplugandplay} which add appropriately chosen constraint-terms to the unconditional score function to steer the process towards generating desirable samples; \duclam{and the Doob’s $h$-transform approach \cite{denker2024deft} which learns a small, auxiliary correction to the unconditional score function.} \duclam{A closely related approach is the Monte Carlo guided diffusion framework (MCGDiff) \cite{cardoso2023monte}, which leverages posterior samples from a linear Bayesian inverse problem obtained via sequential Monte Carlo methods using the singular value decomposition (SVD) to guide the training of a diffusion model. The theoretical foundation of MCGDiff however builds upon the availability of the full SVD of the forward operator, which restricts its applicability to mildly ill-posed problems and moderate-dimensional settings where the SVD can be computed and inverted stably. In contrast, UCoS is designed to operate in a fully \emph{matrix-free} manner, requiring access only to the forward map and its adjoint.}
%While these approaches are convenient as the unconditional score function can be learned independent of the inverse problems, they usually do not have mathematical error bounds as the approximation errors are hard to quantify. Moreover, the information about the forward map has to be incorporated during the sampling process inducing potentially large computational overhead if the forward mapping is expensive to evaluate.

Another line of work seeks to approximate the conditional score function of the posterior distribution directly \cite{saharia2021image,batzolis2021conditional,baldassari2024conditional}. As this approach increases the input dimension of the score function drastically, the training process is more computationally expensive and requires more high quality training samples, which may be restrictive, especially in very high-dimensional problems. 
%In relation to this line of work, we demonstrate that in the context of linear inverse problems the dependence of the conditional score to the measurement data can be simplified and training is not needed in full product space, reducing the computational cost dramatically.
%A second possible approach to avoid the increase in dimensionality is to approximate the conditional score function only after a measurement is made. This is however not suitable for applications in which posterior samples should be generated relatively fast and autonomous after a measurement is made \cite{}.

The theory of infinite-dimensional SDMs has been initiated only very recently. 
\citet{hagemann2023multilevel} modifies the training phase of diffusion models to enable simultaneous training on multiple discretization levels of functions and proves the consistency of their method. \citet{lim2023score} also generalizes the trained model over multiple discretization levels proposing to generate samples with the annealed Langevin algorithm in infinite dimensions. \citet{pidstrigach2023infinitedimensional} was the first to formulate the SDM directly in infinite-dimensional space, demonstrating that the formulation is well-posed and providing theoretical guarantees. Our work is closely connected to \citet{baldassari2024conditional}, where the authors introduce the conditional score in an infinite-dimensional setting. Moreover, they provide a set of conditions to be satisfied to ensure robustness of the generative method and prove that the conditional score can be estimated via a conditional denoising score matching objective in infinite dimensions.



\subsection{Main contribution}
\duclam{This work presents a novel and scalable framework for posterior sampling in high-dimensional and infinite-dimensional inverse problems. The key insight is that posterior samples generated via conditional SDMs can be computed without any evaluation of the forward operator during sampling.}
Instead, without introducing any errors, the computational effort can be shifted to the offline task of training the \emph{unconditional} score of a specific diffusion-like random process. This foundational principle generalizes to other infinite-dimensional diffusion models, beyond the Ornstein-Uhlenbeck process studied here.
More precisely, 
\begin{itemize}
    \item In Theorem \ref{thm:score_id1 infinite dimension} we establish an identity for the conditional score connecting it to a task-dependent unconditional score through affine transformations depending on the forward mapping and the measurement data. The theorem extends this general principle to an infinite-dimensional setting.
    \item In Theorem \ref{thm: stability estimate} we derive an error estimate for the generative process following the ideas suggested by  \citet{pidstrigach2023infinitedimensional}.
    In particular, the error estimate explicitly underscores the contribution of the loss function employed in training.
\end{itemize}
We numerically explore our method in Section \ref{sec: Numerics}, demonstrating that task-dependent training is effective in practice and that online posterior sampling can be performed without evaluations of the forward operator.




%%%% Background section %%%%%

\section{Background}
%\duclam{Lam to write this section}\\

\subsection{Score-based diffusion models in infinite dimensions}
\label{subsec:score-diffusion_in_infdim}
\paragraph{Finite-dimensional score-based diffusion models}
Score-based diffusion models (SDMs) are state-of-the-art machine learning generative models (\cite{song2021scorebasedDiffusion}) that learn a data distribution through gradual denoising of a normally distributed random variable. A diffusion process diffuses an image $X_0\sim p_0$ from $t=0$ to $t=T$ via the SDE (Ornstein-Uhlenbeck process, or OU)
\begin{equation*}
    \dd X_t = -\frac{1}{2}X_t \dd t + \dd W_t,
\end{equation*}
where the marginal densities are denoted by $p_t$, in particular, for large $T$, $p_T$ is close to the  Gaussian density $\mathcal{N}(0,I)$. The marginal densities are then reversed using the backward SDE
\begin{equation*}
    \dd Y_t = \frac{1}{2}Y_t \dd t + \nabla_x\log p_{T-t}(Y_t)\dd t + \dd W_t,
\end{equation*}
such that $Y_T \sim p_0$.

\paragraph{Infinite-dimensional score-based diffusion models}
Let us now review the unconditional SDMs in infinite dimensions proposed by \citet{pidstrigach2023infinitedimensional}. Let $\mu$ be the target distribution, supported on a separable Hilbert space $(H, \langle \cdot, \cdot \rangle_H)$.
Let $\{X^\mu_t\}_{t=0}^T$ stand for the infinite-dimensional diffusion process for a continuous time variable $t\in [0,T]$ satisfying the following infinite-dimensional SDE
\begin{eqnarray}\label{eq: diffusion model; initiated from posterior} 
    dX^{\mu}_t = -\frac{1}{2} X^{\mu}_t dt +  \C^{1/2} dB_t,
    \end{eqnarray}
where $\C$ is a fixed trace class, positive-definite, symmetric covariance operator $\C: H\to H$ and $B_t$ is a Wiener process on $H$ with covariance $tI$,
see Appendix \ref{subsec: Wiener process definition}. 
We assume that the process  $X^{\mu}_t$ is initialized with $\mu$, i.e., $X^\mu \sim \mu$. Notice carefully that we embed the initial condition $\mu$ to the notation $X^\mu_t$, $t\geq 0$, as we later analyze the interplay between different initializations. Here and in what follows, we assume that the initial conditions and the driving process are independent. 
We adopt the formal definition of the score functions by \citet{pidstrigach2023infinitedimensional}:

\begin{definition}\label{def: definition of OU Score function in infinite dimension}
    Define the (weighted, unconditional) score function $s(x,t; \mu)$ for $x \in H$ as
    \begin{eqnarray}\label{eq: definition of OU Score function in infinite dimension}
       s(x, t; \mu) = -\frac{1}{1-e^{-t}}\left(x - e^{-t/2}\E(X^{\mu}_0 |X^{\mu}_t = x)\right).
    \end{eqnarray}
\end{definition}

%[Score for the  process in finite dimension]
\begin{remark}
    Assume that $H = \R^n$ and that the distribution $\mu$ admits a density $p_0$ with respect to the Lebesgue measure.  For any $t \ge 0$ the random variable $X^{\mu}_t$ obtained through SDE in \Eqref{eq: diffusion model; initiated from posterior} has a density $p_t$. The unconditional score $s(x,t; \mu)$ given in \Eqref{eq: definition of OU Score function in infinite dimension} then satisfies $s(x, t; \mu) = \C \nabla_x \log p_t(x)$.
    % \begin{eqnarray*}
    %     s(x, t; \mu) = \C \nabla_x \log p_t(x).
    % \end{eqnarray*}
    In other words, in finite dimensions, Definition \ref{def: definition of OU Score function in infinite dimension} reduces to the common score identified as the log-gradient of the density, scaled by the covariance matrix $\C \in \R^{n\times n}$.
\end{remark}
Now assume that for $T> 0$, 
\begin{equation} \label{eq:expected_square_norm_of_score_is_finite}
    \sup_{t \in [0, T]} \E \norm{s(X_t^\mu, t; \mu)}_H^2 < \infty,
\end{equation}
It is shown in \citet{pidstrigach2023infinitedimensional} that   
given $Y_T^{\mu} \sim \L(X_T^{\mu})$ the following SDE
    \begin{eqnarray}\label{eq: BSDE infdim}
    dY^{\mu}_t = \left(-\frac{1}{2} Y^{\mu}_t- s(Y^{\mu}_t, t; \mu)\right) dt + \C^{1/2} d\bar B_t
\end{eqnarray}
is the time-reversal of \Eqref{eq: diffusion model; initiated from posterior},
where $\bar B_t$ is a different Wiener process on $H$.
The process $Y^{\mu}_t$ is independent of the past increments of $\bar B_t$, but not of the future ones, see \citet{ANDERSON1982-BSDE}.


\paragraph{Training and sampling}
In both finite- and infinite-dimensional diffusion models, the score function is learned by a neural network $s_\theta(\cdot,\cdot;\mu): H \times [0,t] \to H$ such that $s_\theta(x,t; \mu) \approx s(x,t;\mu)$ (where in finite dimension $H=\R^n$ we interpret $s(x,t;\mu)$ as $ \nabla\log p(x)$). 
A common technique to enable empirical learning of the score function is conditional denoising score matching, introduced in \cite{Vincent2011ACB-Loss-function}. This method has been shown to work in the conditional \cite{batzolis2021conditional} and infinite-dimensional setting \cite{baldassari2024conditional}  for the forward SDE given by the OU process. 
Training is performed via stochastic gradient descent by minimizing a score-matching objective (loss function),
\begin{equation}
    \mathrm{SM}(s_\theta) = \E_{t \sim U[\delta, T]} \E_{x \sim \L(X_t^\mu)} \lambda(t)^2\|s_\theta(x,t; \mu) - s(x,t;\mu)\|^2_H,
\end{equation}
over some appropriate class of neural network. Here $\lambda: [\delta,T] \to \R^+$ is a positive weighted function and $0 \le \delta < T$ a truncation that avoids numerical instability for small times. 
In general, the score function $s$ is intractable, such that a denoising score-matching objective is introduced,
\begin{equation*}
    \mathrm{DSM}(s_\theta) = \E_{t \sim U[\delta, T]} \E_{(x, x_0) \sim \L(X_t^\mu, X_0^\mu)} \lambda(t)^2\|s_\theta(x,t; \mu) - (1-e^{-t})^{-\frac{1}{2}}(x - e^{-\frac{t}{2}}x_0)\|^2_H.
\end{equation*}
It can be shown, that the two objectives equal up to a constant that does not depend on the parameter $\theta$. 
After training a score $s_\theta(\cdot,t)$ for all time $t\in [0,T]$, we can then generate samples from $\mu$ by running the backward SDE \Eqref{eq: BSDE infdim} with the trained score $s_\theta$ instead of $s$. The solution of the backward SDE can then be numerically solved using traditional methods such as Euler--Maruyama approximations.

\subsection{Bayesian inverse problems and conditional score}

In this paper, we focus on the setting, where the score corresponding to the reverse drift is conditioned on observations.
We consider a linear  inverse problem
\begin{equation}
	\label{eq:ip}
	y = Ax + \epsilon,
\end{equation}
where $x\in H$ is the unknown in some separable Hilbert space $H$, $y\in \R^m$ stands for the measurement and $A: H \rightarrow \R^m$ is a bounded linear operator. The random noise is modeled by a centered Gaussian distribution $\epsilon \sim \N(0, \Gamma)$, $\Gamma\in \R^{m\times m}$.
We adopt a Bayesian approach to inverse problems  \cite{stuart2010inverse} and assume to have some prior knowledge of the distribution of $x$ before any measurement is made. This knowledge is encoded in a given prior $\mu$, defined as a probability measure on $H$. 
This approach gives rise to a posterior that is absolutely continuous with respect to the prior and its Radon--Nikodym derivative is given by, $\mu$-almost everywhere, 
\begin{eqnarray}
\label{eq:posterior}
     \frac{\dd \mu^y}{\dd\mu}(x)  & = & \frac{1}{Z(y)} \exp\left(-\frac 12 \norm{Ax -y}_\Gamma^2\right), \\ 
     Z(y) & = & \int_H \exp\left(-\frac 12 \norm{Ax -y}_\Gamma^2\right) \mu(\dd x).
     \nonumber
\end{eqnarray}
%. The normalising constant $Z$ is given by
%\begin{equation*}
%	Z(y) = \int \exp\left(-\frac 12 \norm{Ax -y}_\Gamma^2\right) \mu(dx)
%\end{equation*}
In this context, \citet{baldassari2024conditional} defines the conditional infinite-dimensional score as follows:
\begin{definition}\label{def:conditional_score_infiniteD}
    Define the (weighted) conditional score function $s(x,t; \mu^y)$ for $x \in H$ as
\begin{equation*}
%\label{eq:conditional_score_infiniteD}
       s(x, t; \mu^y) = \frac{e^t}{1-e^{t}}\left(x - e^{\frac{-t}{2}}\E(X^{\mu}_0 | Y = y, X^{\mu}_t = x)\right).
\end{equation*}
\end{definition}

\begin{remark} \label{rm:cond_score}
Note the dependency of the conditional score on the measurement $y$. In the same spirit of \Eqref{eq:expected_square_norm_of_score_is_finite}, \citet{baldassari2024conditional} assume the uniform boundedness in time of the expected square norm of $s(x,t;\mu^y)$ for all $y\in\R^m$ to prove the well-posedness of the corresponding time-reversal SDE, given by \Eqref{eq: BSDE infdim} with $\mu $ replaced by $ \mu^y$ and $s(x,t;\mu)$ by $s(x,t;\mu^y)$.  %the score replaced by the one in Def. \ref{def:conditional_score_infiniteD}.
\end{remark}



%%%%% Main theory section %%%%

\section{UCoS: theory and implementation}\label{sec:UCoS}
%

In this section, we establish the theoretical foundation for our method by showing that the conditional score associated with the posterior distribution in the linear inverse problem in \Eqref{eq:ip} can be expressed as an affine transformation of an unconditional score derived from a modified diffusion process. This key observation allows us to shift the complexity of posterior sampling to an offline training phase, where only the prior distribution and forward model are used. 
%
\subsection{UCoS in finite-dimensional space}\label{sec: finite dimensional UCoS}
To build intuition, we begin by illustrating the idea in the finite-dimensional setting $H=\R^n$ \fabian{and assume $\C = I$ for simplicity.}
Consider the posterior $\mu^y$ in \Eqref{eq:posterior} diffused by \Eqref{eq: diffusion model; initiated from posterior}. Suppose the probability densities of the prior $\mu$ and $X_t^{\mu^y}$ are absolutely continuous w.r.t. Lebesgue measure and denote them by $p_0$ and $q_t(\cdot | y)$, respectively. The transitional probability for the multivariate OU process is given by $\rho_t(x | x_0) = {\mathcal N}(e^{-t/2}x_0, (1-e^{-t})\C)$ and thus
\begin{equation}
    \begin{aligned}
    \label{eq:qt_identity}
    &q_t(x|y) =  \int_{\R^n} q_0(z|y) \rho_t(x | z) dz \\
    & = \frac{1}{Z(y)}\int_{\R^n}  p_0(z) \underbrace{\exp\left(-\frac 12 \norm{Az-y}^2_\Gamma\right) \rho_t(x | z)}_{\text{complete square}}dz \\
    & = \frac{\omega_t(x,y)}{Z(y)}\int_{\R^n}  p_0(z) \exp\left(-\frac 12 \norm{z-m_t(x,y)}_{\Sigma_t}^2\right) dz,
    \end{aligned}
\end{equation}
where we completed the square w.r.t. the product of likelihood and transitional probability $\rho_t$. Here, the quantities $\Sigma_t$, $m_t(x,y)$ and $\omega(x,y)$ satisfy
\begin{eqnarray*}
    \Sigma_t &=& \left( \frac{1}{e^t-1}I + A^\top \Gamma^{-1}A\right)^{-1}, \\
    m_t(x,y) & = & \Sigma_t \left(A^\top \Gamma^{-1} y + \frac{x}{e^{t/2} - e^{-t/2}}  \right) \; \text{and} \\
    \omega_t(x,y) & \propto & \exp\left(-\fabian{\frac{1}{2}}\norm{y-e^{t\fabian{/2}} Ax}_{C_t}^2\right),
\end{eqnarray*}
where $C_t = (e^t - 1) A A^\top + \Gamma\in \R^{m\times m}$. Consequently, we have that 
\begin{equation}
    \begin{aligned}
        \label{eq:finite_dim_score_id}
    \nabla_x \log q_t(x|y) 
    =& -\frac12 \nabla_x \norm{y-e^{t/2} Ax}_{C_t}^2  + \nabla_x \log \left((p_0 \ast \N(0,\Sigma_t))(m_t(x,y))\right) \\
    = & - \frac12\nabla_x \norm{y-e^{t/2} Ax}_{C_t}^2 + (\nabla_x m_t(x,y))^\top  \left[\nabla_x \log \left(p_0 \ast \N(0,\Sigma_t)\right)\right](m_t(x,y)).
    \end{aligned}
\end{equation}
The identity in \Eqref{eq:qt_identity} can be formally interpreted as expressing $q_t$ for a given time %$†$ 
as a mixture of the prior $p_0$ and a Gaussian distribution with time-dependent covariance $\Sigma_t$ modulo transformations with $\omega_t$ and $m_t$. What is more, the conditional score $s(\cdot; \mu^y)$ corresponding to $q_t$ in \Eqref{eq:finite_dim_score_id} can be expressed as a transformation of the score $\tilde s$ of this mixture, $\tilde s(x; t; \mu) :=\Sigma_t\nabla_x \log \left(p_0 \ast \N(0,\Sigma_t)\right)$. This idea enables us to approximate the posterior score offline up to affine transformations. More precisely, we have the following identity:
\begin{equation}\label{eq:finite_dim_score_id 2}
\begin{aligned}
     s(x, t; \mu^y) &= \lambda(t) \left(  \Tilde{s}\left(m_t(x, y\right), t; \mu) + (e^t-1) A^\top C_t^{-1}(y-e^{t/2}Ax)\right) \\
    & =\lambda(t) \left(  \Tilde{s}\left(m_t(x, y\right), t; \mu) + m_t(x, y) - e^{t/2}x \right),
\end{aligned}
\end{equation}
with
\begin{eqnarray*}
   \lambda(t) =  (e^{t/2} - e^{-t/2})^{-1}.
\end{eqnarray*}
%Lemma \ref{lem: writing m_t for no forward} was used to simplify.


\begin{remark}
The principled idea of transforming from task-dependent unconditional score to the conditional score outlined above generalizes to other linear diffusion models beyond the specific Ornstein-Uhlenbeck process \Eqref{eq: diffusion model; initiated from posterior} studied here. This raises the important question of how to design an efficient underlying diffusion model to balance the computational effort further in a desirable way, e.g. by temporal or spectral weighting of the diffusion. This consideration is beyond the scope of this paper.
\end{remark}

The result of \Eqref{eq:finite_dim_score_id 2} motivates a new way to approximate the score function in a task-dependent manner. Notice first that the term $m_t(x,y)$ depends on the forward map $A$ in a non-trivial and time-dependent way. Therefore, even if training of $\tilde s$ can be performed offline and utilized as an approximation in \Eqref{eq:finite_dim_score_id 2}, one would still need to evaluate $A$ during the sample generation. In what follows, we propose an approximation scheme for the conditional score that circumvents this issue.
Notice that the dependence of $m_t$ on $A$ is mostly through multiplication with $\Sigma_t$. Hence we may learn the function $r$ given by
\begin{equation*}
    r(\cdot, t; \mu) := \tilde{s}(\Sigma_t \cdot, t; \mu) + \Sigma_t \cdot
\end{equation*}
instead of $\tilde s$. By substituting $r$ into \Eqref{eq:finite_dim_score_id 2}, we obtain the final identity
\begin{equation*}\label{eq:finite_dim_score_id 3}
    s(x, t; \mu^y) = \lambda(t) \left(  r(\xi_t(x;y)- e^{t/2}x \right)
\end{equation*}
with $\xi$ given by
\begin{equation*}	
	\xi_t(x,y) =  A^\top \Gamma^{-1} y + \lambda(t)x.
\end{equation*}
We can compute and save the term $ A^\top \Gamma^{-1} y$ to memory using a single adjoint evaluation. For each time step $t>0$, we add the time-dependent second term to obtain $\xi_t$ without the need for any further forward evaluations during sampling. 

\begin{remark}
    Assume that the prior $p_0 = \N(0, S_0)$ is Gaussian. Then $r$ is given by
    \begin{equation*}
        r (\zeta, t; \mu) =  \left(\frac{1}{e^t-1}\C^{-1} + A^\top \Gamma^{-1}A + S_0^{-1} \right)^{-1} \C^{-1}\zeta
    \end{equation*}
and is a bounded linear mapping for any $t>0$. In particular, $r :\R^n \times \R \to \R^n$ depends only on a spatial variable $\zeta$ and a temporal variable $t$ but \textbf{not} on the measurement $y$. \fabian{This is in stark contrast to the conditional approach \cite{baldassari2024conditional}, where the conditional score 
\lam{$$s: (x,y,t) \mapsto s(x,t; \mu^y) : \R^n \times \R^m \times \R \to \R^n$$ 
(see Definition \ref{def:conditional_score_infiniteD}) depends on \textbf{both} spatial variables $x$ and $y$. The added dependence on $y$ increases the input dimensionality from $n+1$ to $n+m+1$, which substantially raises the required sample complexity and the size of the neural network required to achieve a given error tolerance.}}
\end{remark}
%In the case of Gaussian prior, the conditional score simplifies to:
%\begin{eqnarray*}
%    s(x, t; \mu^y) = -\left( e^{-t} \left(A^\top \Gamma^{-1} A + S_0^{-1}\right)^{-1} +(1-e^{-t} ) I\right)^{-1} \left(x-e^{-t/2} \left(A^\top \Gamma^{-1} A + S_0^{-1}\right)^{-1} A^\top \Gamma^{-1}y\right).
%\end{eqnarray*}

\subsection{Theoretical foundations of UCoS in function space}
We will now make the transformations above precise in infinite-dimensional separable Hilbert space $H$. First, we define a random process $\{\widetilde{X}^{\mu}_t\}_{t=0}^T$ such that its distribution is given by the mixture model
\begin{equation}\label{eq: adapted forward model}
    \widetilde X^{\mu}_t = \widetilde X^{\mu}_0 + Z_t, \quad \widetilde X^{\mu}_0 \sim \mu \; \text{and} \; Z_t \sim \N(0, \Sigma_t)
\end{equation}
with $\Sigma_t$ given  by
\begin{eqnarray}\label{eq: formal definition of Sigma_t}
   \Sigma_t = (e^t -1)\C - (e^t-1)^2\C A^* C_t^{-1} A \C
\end{eqnarray}
and
\begin{equation}
    \label{eq:Ct}
	C_t = (e^{t}-1) A \C A^* + \Gamma \in \R^{m \times m}.
\end{equation}
\fabian{Notice that $\Sigma_t$ coincides with its formulation in Section \ref{sec: finite dimensional UCoS} in finite-dimensional case, see Lemma \ref{lem: relation between CM spaces of U and Vt}.}
The next lemma ensures that $\Sigma_t$ is a valid covariance operator.
\begin{lemma}\label{lem: Sigma_t is covariance operator}
For $t>0$, the operator $\Sigma_t : H \to H$ is trace class, self-adjoint and positive definite.
\end{lemma}

\begin{remark}
While the mixture in \Eqref{eq: adapted forward model} is well-defined for any $t>0$, we observe that the process $\widetilde X_t^{\mu}$ is no longer diffusive in the sense that it may stay dependent of $\widetilde X_0^{\mu}$ as $t$ grows. Indeed, we notice the dual behavior from two cases: 
\begin{itemize}
    \item if $A = I$, we have $\Sigma_t = \C (\C + (e^t-1)^{-1} \Gamma)^{-1} \Gamma$.
    \item for $A=0$ we have $\Sigma_t = (e^t-1)\C$. 
\end{itemize} 
This indicates different asymptotics for the variance of $\widetilde X_t^{\mu}$ depending on the singular values of $A$.
In order to improve intuition on the correlation structure of $\tilde X_t$, we depict the diagonal of $\Sigma_t$ for different values of $t$ and the inverse problems related to inpainting in Figure \ref{fig: Sigmat inpainting} and CT imaging in Figure \ref{fig: Sigmat ct}. A detailed description of these problems can be found in Section \ref{sec: Numerics}. We further note studying an SDE corresponding to $\widetilde X_t^{\mu}$ or its time-reversal is beyond the scope of this paper.
\begin{figure}[ht]
    \centering
    \begin{minipage}{0.24\textwidth}
        \centering
        \includegraphics[width=\linewidth]{Images/Sigma_t_structure/Sig_tinpaint0.05.png}
    \end{minipage}%
    \begin{minipage}{0.24\textwidth}
        \centering
        \includegraphics[width=\linewidth]{Images/Sigma_t_structure/Sig_tinpaint0.1.png}
    \end{minipage}%
    \begin{minipage}{0.24\textwidth}
        \centering
        \includegraphics[width=\linewidth]{Images/Sigma_t_structure/Sig_tinpaint0.5.png}
    \end{minipage}%
    \begin{minipage}{0.24\textwidth}
        \centering
        \includegraphics[width=\linewidth]{Images/Sigma_t_structure/Sig_tinpaint0.9.png}
    \end{minipage}
    \caption{Diagonal of $\Sigma_t$ for inpainting problem. Time-steps from left to right: $t = 0.05, t = 0.1, t = 0.5$ and $t = 0.9$.}
    \label{fig: Sigmat inpainting}
\end{figure}
\begin{figure}[ht]
    \centering
    \begin{minipage}{0.24\textwidth}
        \centering
        \includegraphics[width=\linewidth]{Images/Sigma_t_structure/Sig_tct0.05.png}
    \end{minipage}%
    \begin{minipage}{0.24\textwidth}
        \centering
        \includegraphics[width=\linewidth]{Images/Sigma_t_structure/Sig_tct0.1.png}
    \end{minipage}%
    \begin{minipage}{0.24\textwidth}
        \centering
        \includegraphics[width=\linewidth]{Images/Sigma_t_structure/Sig_tct0.5.png}
    \end{minipage}%
    \begin{minipage}{0.24\textwidth}
        \centering
        \includegraphics[width=\linewidth]{Images/Sigma_t_structure/Sig_tct0.9.png}
    \end{minipage}
    \caption{Diagonal of $\Sigma_t$ for CT imaging problem. Time-steps from left to right: $t = 0.05, t = 0.1, t = 0.5$ and $t = 0.9$.}
    \label{fig: Sigmat ct}
\end{figure}
\end{remark}




    
\begin{definition}
We define the (weighted, unconditional) score function of the process in \Eqref{eq: adapted forward model} by
    \begin{eqnarray}\label{eq: posterior score in infinite dimensional space}
    \tilde{s}(z, t; \mu) = - \left(z -\E(\widetilde X^{\mu}_0 |\widetilde X^{\mu}_t = z)\right).  
\end{eqnarray}
\end{definition}
\begin{lemma}\label{lem: Score function of posterior process in infinite dimension}
Assume that $H = \R^n$ and our prior measure $\mu$ admits a density given by $p_0$. Moreover, let $\widetilde X_t^\mu$ be defined by \Eqref{eq: adapted forward model}.
Then for every $t \ge 0$, the random variable $\widetilde X^{\mu}_t$ admits a density $\tilde p_t(x)$ satisfying
\begin{eqnarray*}
    \tilde s(z, t; \mu) = \Sigma_t \nabla_z \log \tilde p_t(z).
\end{eqnarray*}
\end{lemma}
 Let us fix some $t>0$ and define the operator $R_t: H\to H$ through
\begin{equation}
    \label{eq:Rt}
    R_t := (e^t-1) I + (e^t-1)^2 \C A^* C_t^{-1} A.
\end{equation}
By Lemma \ref{lem: Useful equality that show that Sigma is still p.d.}, $R_t$ is bijective.
Define the affine transformation $r$ of $\tilde s$ by
\begin{equation}\label{eq: r equation}
    r(\cdot, t; \mu) := \tilde{s}(R_t \cdot, t; \mu) + R_t\cdot,
\end{equation}

Let us now proceed to the main result of this section.
\begin{theorem}\label{thm:score_id1 infinite dimension}
Let $H_{\C}$ be the Cameron-Martin space of $\C$ (see Appendix \ref{Appendix:prob_measures}) and assume that the prior satisfies $\mu(H_{\C}) = 1$.
Then the for conditional score function $s(x,t;\mu^y)$ it holds that
\begin{equation}
    \label{eq:conditional_score_id}
    s(x, t; \mu^y) = \lambda(t) \left(  r\left(\xi_t(x, y)\right)- e^{t/2}x \right)
\end{equation}
for $(x,y) \in H \times \R^m$ a.e. in $\L(X_t^\mu, Y)$ and $t>0$, where
\begin{equation}
	\label{eq:lambda}
	\lambda(t) = (e^{t/2} - e^{-t/2})^{-1}, \qquad \xi_t(x,y) =  \C A^* \Gamma^{-1} y + \lambda(t) x.
\end{equation}
\end{theorem}


\paragraph{Special case of Gaussian prior}
Let us now study the score $\tilde s$ corresponding to the special case of a Gaussian prior measure. This enables us to derive an explicit formula which can give insights, in particular, regarding the regularity of the score function.

\begin{lemma}\label{lem: score function is compact}
       Let $\mu = \N(0, S_0)$ be a Gaussian measure and suppose the covariance operator $\widetilde{C}$ satisfies the assumptions of Theorem \ref{thm:score_id1 infinite dimension}. Then there exists a covariance operator $\C$ on $H$ such that $\Sigma_t (\Sigma_t + S_0)^{-1}: H\to H$ can be well-defined as a linear and bounded operator for $t>0$.
       Moreover, in that case
        \begin{eqnarray*}
        \tilde s(z, t; \mu) & = & -\Sigma_t (\Sigma_t + S_0)^{-1}z.
\end{eqnarray*}
\end{lemma}




\subsection{Matrix-free implementation}
\label{sec:matrix_free_implementation}
Building on the theoretical framework established in the previous section, we now describe how UCoS can be implemented efficiently in practice. A key strength of our approach is that it remains entirely matrix-free: it requires only the ability to evaluate the forward operator \(A\) and its adjoint \(A^*\) without ever forming or inverting these operators explicitly. \fabian{In a discretized setting, this corresponds to avoiding computing and saving the matrix representing $A$, which can be expensive in terms of computational time and memory for large scale problems.} This property enables UCoS to scale to high-dimensional and ill-posed inverse problems where standard representations (e.g., via SVD or matrix factorizations) become computationally infeasible. 
In what follows, we outline the two-phase implementation pipeline: an offline training phase in which a task-dependent score approximation is learned, and an online sampling phase where posterior samples are generated using only the trained network and a precomputed observation-dependent shift.
%
%
\subsubsection{Offline phase: learning the task-dependent score}


The next lemma adapts the training procedure discussed in Section \ref{subsec:score-diffusion_in_infdim} to the setting of forward process $\widetilde X^{\mu}$. Let us define continuous-time score matching objective (loss function) and the denoising score-matching objective 
\begin{equation*}\label{eq:loss-objective}
     \widetilde{\mathrm{SM}}(r_\theta) := \E\left[\lambda(t)^2\norm{\Tilde{s}(\tilde x_t, t; \mu) + \tilde x_t -r_\theta(R_t^{-1} \tilde x_t, t; \mu)}_H^2\right],
\end{equation*}
    \begin{equation*}
        \widetilde{\mathrm{DSM}}(r_\theta) := \E \left[\lambda(t)^2 \norm{r_\theta(R_t^{-1} \tilde x_t, t; \mu) - \tilde x_0}_H^2\right].
    \end{equation*}
with the expectations taken w.r.t. $ t \sim U[\delta, T], (\tilde x_0, \tilde x_t) \sim \L(\widetilde{X}^{\mu}_0, \widetilde{X}_t^{\mu})$.
We make an assumption regarding the boundedness of the second moment of the conditional score uniformly in time following \citet{baldassari2024conditional}.
\begin{assumption}\label{ass: stability estimate; no forward process}
       The prior $\mu$ has bounded second moment, $\E_{X\sim\mu}\|X\|_H^2 < \infty$, and
    \begin{eqnarray*}
    \sup_{t \in [\delta, T]} \E_{y \sim \pi_y} \E_{x_t \sim \L(X_t^{\mu^y})} \norm{ s(x_t, t; \mu^y)}_H^2 < \infty.
\end{eqnarray*}

\end{assumption}
Notice that Assumption \ref{ass: stability estimate; no forward process} does not require us to sample from $Y \sim \pi_y$ or to compute the transform $m_t(x, y)$. See Lemma \ref{lem: training loss} in the appendix for detailed justification.
\begin{lemma}\label{lem: conditional score matching}
    Let Assumption \ref{ass: stability estimate; no forward process} hold. It follows that
    \begin{eqnarray*}
        \widetilde{\mathrm{SM}}(r_\theta)
        = \widetilde{\mathrm{DSM}}(r_\theta) + V,
    \end{eqnarray*}
    where the constant $V < \infty$ is independent of $\theta$.
\end{lemma}
The result follows from repeating the arguments in \citet{baldassari2024conditional}. 
    Truncation $t\geq \delta$ guarantees that complications with singularity at $t\to 0$ are avoided, see Appendix \ref{sec: proof of conditional score matching} for details.


By the definition of the process $\widetilde{X}_t^{\mu}$, we have $\widetilde{X}_t^{\mu} = \widetilde{X}^{\mu}_0 + Z_t$ in distribution, with $Z_t \sim N(0, \Sigma_t)$. 
Using this identity, we observe that $R_t^{-1} \tilde x_t$ conditioned on $\tilde x_0$ is Gaussian with covariance
$$ \text{Cov}(R_t^{-1}\widetilde{X}_t^{\mu} \; | \; \widetilde{X}^{\mu}_0) = R_t^{-1} \Sigma_t \left(R_t^{-1}\right)^* = \frac{1}{e^t-1}{\mathcal C} + {\mathcal C} A^* \Gamma^{-1}A {\mathcal C}$$
by the identity $$R_t^{-1} = \frac{1}{e^t-1}I + {\mathcal C} A^* \Gamma^{-1}A,$$
followed from Lemma \ref{lem: Useful equality that show that Sigma is still p.d.} and Lemma \ref{lem: relation between CM spaces of U and Vt} (i).
Now we can directly generate the desired quantity by observing that
$$R_t^{-1} \tilde x_t = \left(\frac{1}{e^t-1}I + {\mathcal C} A^* \Gamma^{-1}A\right) \tilde x_0 + \frac{1}{\sqrt{e^t-1}}{\mathcal C}^{1/2} z_1 + {\mathcal C} A^* \Gamma^{-1/2} z_2 $$
in distribution, where $(\tilde x_0, \tilde x_t) \sim {\mathcal L}(\widetilde{X}^{\mu}_0, \widetilde{X}_t^{\mu})$ and $z_1, z_2 \sim N(0, I)$ i.i.d.
In consequence, the training phase requires application of the forward operator $A$, its adjoint $A^*$, the covariance of the noising process ${\mathcal C}$ and the precision matrix of the observational noise $\Gamma^{-1}$. We note that $\Gamma$ in the case of degenerative noise distribution, $\Gamma^{-1}$ can be replaced by the pseudoinverse.
%Thus, the offline phase requires only access to the operators \(A, A^*, \mathcal{C}, \Gamma^{-1}\), all of which can be implemented matrix-free. In degenerate noise settings, \(\Gamma^{-1}\) may be interpreted as a pseudoinverse.
%
\subsubsection{Online phase: posterior sampling without forward map evaluations}
The offline training, detailed above, will yield an (e.g. FNO) approximation $r_\theta \approx r$ independent of any measurement data. Once this is done, sampling from the posterior is performed by simulating a reverse-time diffusion process using the transformed score $s$, defined as (taking into account Theorem \ref{thm:score_id1 infinite dimension})
$$ s(x,y; \mu^y) = \lambda(t) \left(r_\theta(\xi_t(x,y), t; \mu) - e^{-\frac t2} x\right), $$
where 
\begin{eqnarray*}
    \xi_t(x,y) = {\mathcal C}A^* \Gamma^{-1}y + \lambda(t) x.
\end{eqnarray*}
Note that after collecting a new measurement vector $y$, the term ${\mathcal C}A^* \Gamma^{-1}y$ is computed once before the generative process and remains constant throughout it. In consequence, the generative process does not require evaluation of any operators in addition to $r_\theta$.


%%%%%% Convergence analysis section %%%%
\section{Convergence analysis}\label{sec: Estimating numerical errors}

The aim of this section is to conduct rigorously the theoretical convergence of the proposed method. More precisely, we establish a quantitative bound of the error term that indicates how far the samples generated by UCoS lie from the true posterior target measure. In particular, we quantify how this error term depends on different types of numerical approximations. Let us fix $T >0$ and $\delta \in (0, T)$. 
Consider a partition $\{\delta = t_1 \le \cdots \le t_n = T \}$ of $[\delta, T]$ with mesh size $\Delta t := \min\{t_i - t_{i+1} | 0 \le i \le N-1\}$ and define $\low{t}:= \max \left\{t_i | 1 \le i \le n, \; t_i \le t\right\}$. 
For clarity, we reverse time in \Eqref{eq: BSDE infdim} and let $W_t^{\mu^y} = Y_{T-t}^{\mu^y}$ for $y$ $\pi_y$-a.e. be the ideal solution satisfying
\begin{equation}\label{reverse time eq}
    dW_t^{\mu^y} = \left(-\frac{1}{2} W_t^{\mu^y} - s(W_t^{\mu^y}, T-t; \mu^y) \right) dt + \C^{1/2} dB_t
\end{equation}
initiated at $w_0 = (1-e^{-T}) z + e^{- T/2} x_0$ for $z \sim \N(0, \C)$ and $x_0 \sim \mu^y$.
Observe that \Eqref{reverse time eq} is a `standard' SDE where $W_t^{\mu^y}$ is independent of the future increments of $B_t$. Moreover, let $v_t$ correspond the numerical approximation to $W_t^{\mu^y}$ satisfying
\begin{equation*}
\begin{aligned}
    d  V_t^{\mu^y} = & -\frac{1}{2}  V_\low{t}^{\mu^y} - \lambda(T-\low{t})\left[ r_\theta\left(\xi_{T-\low t}(V_{\low{t}}^{\mu^y}, y), T-\low{t}; \mu\right) \right. \left. - e^{(T-\low{t})/2} V_{\low{t}}^{\mu^y} \right] dt +\C^{1/2} dB_t
\end{aligned}
\end{equation*} 
and initiated at $v_0 =  (1-e^{- T}) z$. Here, the two stochastic processes share the same Wiener process $B_t$ and initialization $z$. In what follows, we consider a discrete-time loss given by
\begin{equation}
\label{eq:discrete_loss}
\begin{aligned}
    &\varepsilon_{LOSS} =  \E_{t \sim U[\delta, T]} \E_{\tilde x_{\low t} \sim  \L(\widetilde{X}_{\low t}^{\mu})} \lambda(\low t)^2\norm{\Tilde{s}(\tilde x_{\low t}, \low t; \mu) + \tilde x_{\low t} -r_\theta(R_{\low t}^{-1} \tilde x_{\low t}, \low t; \mu)}_H^2,
\end{aligned}
\end{equation}
We are now ready to state the main result of this section.
\begin{theorem}\label{thm: stability estimate}
Let Assumption \ref{ass: stability estimate; no forward process} and the assumptions of Theorem \ref{thm:score_id1 infinite dimension} hold. Assume further that $s(\cdot,t;\mu)$ and $r_\theta(\cdot,t;\mu)$ are Lipschitz continuous with Lipschitz constant $L_s(\cdot) \in L^2([\delta,T])$. %and assume that some measurement $y$ from the random variable $Y \sim \pi_y$ is drawn randomly, where $\pi_y$ is the marginal distribution of $y$ in the inverse problem \Eqref{eq:ip}.
Then, %we have the following bound 
\begin{eqnarray*}
    \E_{y \sim \pi_y} \E_{(w_T, v_{T-\delta}) \sim \L( W_T^{\mu^y},V_{T-\delta}^{\mu^y})} {\norm{w_T -  v_{T-\delta}}_H^2} 
    \le M\left( \varepsilon_{NUM} + \varepsilon_{LOSS} + \varepsilon_{INIT} + \delta \right)
\end{eqnarray*}
where $\varepsilon_{LOSS}$ is given in \Eqref{eq:discrete_loss}, and 
\begin{equation*}
    %\varepsilon_{TRUNC} \le  {\mathcal O}(\delta), \; 
    \varepsilon_{INIT} \le e^{-T} \E_{X\sim\mu}\|X\|_H^2, \;
    \varepsilon_{NUM} \le {\mathcal O}(\Delta t),
\end{equation*}
%\begin{eqnarray*}
%     \varepsilon_{LOSS} := \int_\delta^T \E_{\tilde x_0 \sim \L(\widetilde{X}^{\mu}_0)} \E_{\tilde x_\low{t} \sim \L(\widetilde{X}_\low{t}^{\mu} | \widetilde{X}_0^{\mu})} \left[\lambda(\low{t})^2\norm{\Tilde{s}(\tilde x_\low{t}, \low{t}; \mu) + \tilde x_\low{t} -r_\theta(T_{T-\low{t}}^{-1} \tilde x_\low{t}, \low{t}; \mu)}_H^2\right],
%\end{eqnarray*}
and the constant $M$ depends on the quantities
\begin{eqnarray*}
    \E_{t\sim U[0,T-\delta]}\E_{y \sim \pi_y}\E_{w_t \sim \L(w_t)} \norm{s(w_t, T-t; \mu^y)}_H^2,  \E_{X\sim\mu}\|X\|_H^2, \; \int_\delta^{T} L_s^2(\tau) d \tau, \; \tr_H (\C)\; \text{ and } \; T.
\end{eqnarray*}
\end{theorem}
%Note that by assumption \ref{ass: stability estimate; no forward process}, $\pi_y$ almost surely
%\begin{eqnarray*}
%    \sup_{t \in [\delta, T]} \E_{x_t \sim \L(X_t^{\mu^y})} \norm{ s(x_t, t; \mu^y)}_H^2 < \infty.
%\end{eqnarray*}
Note that the Lipschitz continuity assumption of the drift is a natural setup for the backward SDE to retain the uniqueness of a (strong) solution, see \citet{pidstrigach2023infinitedimensional}. It also follows from there that for $\pi_y$-almost every $y \in \R^m$ the random variable $W_T^{\mu^y}$ is distributed according to the posterior $\mu^y$.





%%% Numeric section %%%%

\section{Numerical experiments}\label{sec: Numerics}


% It would now be standard to train a neural network $r_\theta(x, t; \mu)$, using the loss function
% \begin{eqnarray*}
%     \E_{t \sim [\delta, T]} \E_{\tilde x_0 \sim \L(\widetilde{X}^{\mu}_0)} \E_{\tilde x_t \sim \L(\widetilde{X}_t^{\mu} | \widetilde{X}_0^{\mu})}  \norm{\frac{\tilde{\mathfrak{s}}_\theta(\tilde x_t, t; \mu) + \tilde x_t - \tilde x_0}{e^{t/2}-e^{-t/2}}}^2.
% \end{eqnarray*}
% Notice however that $\Sigma_t$ blows up exponentially fast if $A$ is not of full rank, hence sampling from $\widetilde{X}_t^{\mu}$ may cause overflow for large $t$. To combat this, we introduce a scaled version 
% \begin{eqnarray*}
%     S_t := \frac{1}{e^t-1} \Sigma_t = \C - (e^t-1)\C A^* C_t^{-1} A \C. 
% \end{eqnarray*}
% and corresponding process
% \begin{eqnarray*}
%     (\widetilde{Z}_t^{\mu})_{0 < t \le T} := \left( \frac{\widetilde{X}_{t}^{\mu}}{\sqrt{e^{t}-1}} \right)_{0 < t \le T}  
% \end{eqnarray*}
% such that
% \begin{eqnarray*}
%     \widetilde{Z}_t^{\mu} |  \widetilde{X}_{0}^{\mu} \sim \N\left(\frac{1}{\sqrt{e^t-1}} \widetilde{X}_{0}^{\mu}, S_t \right).
% \end{eqnarray*}
% In order to obtain an approximation
% \begin{eqnarray*}
%     \tilde{\mathfrak{s}}_\theta\left(
%     x / \sqrt{e^t-1} , t\right) \approx \frac{\tilde{s}(x, t; \mu)}{\sqrt{e^t-1}}
% \end{eqnarray*}
% we choose the loss function
% \begin{eqnarray*}
%     \varepsilon_{LOSS}(\theta) =\E_{t \sim [\delta, T]} \lambda(t) \E_{\tilde x_0 \sim \L(\widetilde{X}^{\mu}_0)} \E_{\tilde z_t \sim \L(\widetilde{Z}_t^{\mu} | \widetilde{X}_0^{\mu})} \norm{\tilde{\mathfrak{s}}_\theta(\tilde z_t, t; \mu) + \tilde z_t - \frac{\tilde x_0}{\sqrt{e^t-1}} ) }^2,
% \end{eqnarray*}
% which does not require us to compute the process $\widetilde{X}_t^{\mu}$, which may explode for large t. We define the time dependent weight $\lambda$ by
% \begin{eqnarray*}
%     {\lambda(t)} = \left(\frac{e^{t/2}}{\sqrt{e^t-1}} \right)^2 = \frac{e^t}{e^t-1} = \frac{1}{1-e^{-t}}.
% \end{eqnarray*}





% \begin{lemma}
%     Note that \cite{Vincent2011ACB-Loss-function} there exists some (time-dependent) constant $V_t$ such that
% \begin{eqnarray*}
%     \varepsilon_{LOSS}(\theta)= \E_{t \sim [\delta, T]} \lambda(t)\E_{Z_t |\widetilde{X}^{\mu}_0} \E_{\widetilde{X}^{\mu}_0}   \norm{\tilde{\mathfrak s}_\theta(Z_t, t; \mu) - \frac{\tilde{s}(Z_t,t; \mu)}{\sqrt{e^t-1}}}^2 + V_t.
% \end{eqnarray*}
% \end{lemma}
% \begin{proof}
%     mostly todo, compare Proposition 4 in \cite{baldassari2024conditional}
% \end{proof}

% \begin{remark}
%     If $\Sigma_t$ is invertible, one could find it tempting to approximate the function $\tilde s(x,t; \mu) $ as $\tilde s := \Sigma_t^{-1} \left( \tilde s\right)$ instead of $\tilde s(x,t; \mu) $. We find this to behave very problematically in our numerical experiments, to the point where visual quality of generated samples was not acceptable.
% \end{remark}

% \begin{eqnarray*}
%     s(x, t; \mu^y) & = & \frac {1}{(e^{t/2} - e^{-t/2})}\left(  \Sigma_t  \Tilde{s}(m_t, t; \mu) + m_t - e^{t/2}x \right)
% \end{eqnarray*}
% with again with the intention of finding a numerical function approximation $\mathfrak s_\theta$ such that evaluating the forward operator is no longer necessary in each iteration during sample generation. Notice that 
% the influence of the forward operator is only through $m_t$, which appears as argument in the score function and addition to the score.

% \begin{eqnarray*}
%     m_t = \Sigma_t\C ^{-1}  \left(\C A^* \Gamma^{-1} y + \frac{1}{e^{t/2} - e^{-t/2}}  x \right).
% \end{eqnarray*}


We showcase our method in the context of inverse problems related to a toy inpainting example, computerized tomography (CT) and a deblurring problem. Details of the numerical implementation are described in Appendix \ref{section: Details of NN} and enlarged depictions of posterior samples can be found in Appendix \ref{sec: additional samples}.

\fabian{To enable a principled comparison, we examine the accuracy of posterior sampling using score approximations trained with architectures that offer comparable runtimes. 
Methods that avoid forward evaluations during sampling, such as UCoS and conditional method, exhibit similar runtimes. However, we find that the performance of the conditional method is hindered by a more complex score approximation, especially when neural operators have very few parameters.
In contrast, the computational cost of posterior sampling with unconditional methods increases due to the need for repeated forward evaluations during the sampling process. When comparing unconditional methods to UCoS, deviations from the true posterior are primarily attributed to the approximative nature of the conditioning on measurement information, rather than inaccuracies in the score approximation. This is because the dimensionality of the score approximation is the same for both UCoS and unconditional methods.}

Following \cite{baldassari2024conditional}, we utilize a Fourier neural operator (FNO) \cite{li2021fourier} to parameterise $r_\theta$ and perform training outlined in Section \ref{subsec:score-diffusion_in_infdim} on an Nvidia A100 GPU with 80 GiB. We compare UCoS to four different methods. First, we consider one of the first conditional score approximation considered in \cite{jalal2021robustannealedlangevian, feng2023scorebased} and refer to it by SDE ALD.
Additionally, we implement Diffusion Posterior Sampling (DPS) \cite{chung2023diffusion}, projection based methods for general inverse problems (Proj) \cite{Dey_2024} and the conditional score approach (Conditional) \cite{baldassari2024conditional}. 
Each score approximation is given explicitly in Appendix \ref{section: Details of NN}. 
The numerical implementation can be found at 
\ExplSyntaxOn
\bool_if:NTF \accepted
    {\url{https://github.com/FabianSBD/SBD-task-dependent/tree/UCoS}.}
    {\url{https://github.com/FabianSBD/SBD-task-dependent/tree/UCoS}.}
\ExplSyntaxOff
%\url{https://anonymous.4open.science/r/SBD-task-dependent-E5F5}
\paragraph{Inpainting} In this paragraph, the prior is modeled by a Gaussian distribution with squared exponential covariance structure. This allows us to compute all score functions in closed form without any learning. Uniquely to this example, we are able to sample from the (Gaussian) posterior. 
The forward mapping $A: \R^n \rightarrow \R^n$ is diagonal and with diagonal entries one in the center of the image and zero elsewhere. The information on the center pixels is lost and in the inverse problem we seek to recover the full image. The image resolution is $64 \times 64$ pixels.

Figure \ref{fig:Inpainting posterior samples} and Table \ref{tab: inpainting} illustrate that the posterior ensemble ($N=1000$) generated by UCoS (our method) and conditional method generate samples that are consistent with the true posterior distribution. Since no score needs to be learned, UCoS and Conditional perform very similarly.
\begin{figure}[h]
    \centering
    % First column: Two vertically stacked images
    \begin{subfigure}[c]{0.2\textwidth}
        \centering
        truth
        \includegraphics[width=\textwidth]{Images/Inpainting/truth.jpg} \\ % First image
        \vspace{5pt} % Optional spacing between the two images
        measurement
        \includegraphics[width=\textwidth]{Images/Inpainting/measurement.png}   % Second image
    \end{subfigure}
    % Second column: Single image
    \begin{subfigure}[c]{0.6\textwidth}
        \centering
        \hspace{1.5cm} samples \hspace{3.2cm} bias \hspace{5mm} std
        \includegraphics[width=\textwidth]{Images/Inpainting/summary325.png}
    \end{subfigure}   
    \begin{subfigure}[c]{0.17\textwidth}
        %\phantom{.} 
        %\\
        SDE ALD \\
        \vspace{6mm}
        \\
        DPS \\
        \vspace{6mm}
        \\
        Proj \\
        \vspace{6mm}
        \\
        Conditional \\
        \vspace{6mm}
        \\
        UCoS \\
        \vspace{2mm}
        \\
        true posterior
    \end{subfigure}
   \caption{Summary of posterior samples for inpainting problem. On the left column, the top image is ground truth and the bottom is measurement data. On the right column, 5 posterior samples with true score corresponding to a Gaussian prior and their bias and standard deviation, with methods used from top to bottom: SDE ALD, DPS, Proj, Conditional, UCoS and true posterior.}
    \label{fig:Inpainting posterior samples}
\end{figure}
\begin{table}
    \centering
    \begin{tabular}{c c c c }
         \centering
             method & bias& std & time\\
             \hline
             SDE ALD& 0.1341 & 0.1419 & 202\\  
             DPS & 0.1050 & 0.1242 & 467 \\
             Proj & 0.2008 & 0.1772 & 254 \\
             Conditional & 0.1011 & 0.1238 & 198 \\
             UCoS & 0.1003 & 0.1226 & 197 \\
             True & 0.1013 & 0.1304 & <1 \\
             \hline
        \end{tabular}
    \caption{Summary statistics of all methods for the inpainting problem. Bias and std are computed on 1000 generated posterior samples and averaged over the number of pixels. Computational time is given in seconds.}
    \label{tab: inpainting}
\end{table}


\paragraph{CT imaging}
Here, the forward mapping $A$ models a sparse-view imaging setting with a 45-degree angle of view with 256 equiangular directions. The detector is assumed to have 256 apertures spanning over the width of imaging area (i.e., 256 parallel line integrals per angle are measured) and, consequently, the problem dimensions are given by $m = n = 256^2$. Moreover, the measurement data is corrupted by additive Gaussian noise with signal--to--noise ratio of 20 dB.

Our training data is the Lung Image Database Consortium image collection (LIDC--IRI)--dataset \cite{armato2015lidc}, containing $>200.000$ 2D slices of resolution $512 \times 512$, we rescale them to $256 \times 256$. 

We analyze the influence of the complexity of the FNO approximation on the posterior samples. We employ an architecture with 4 layers while varying the number of nodes uniformly across all layers. For the full $256 \times 256$ resolution, we analyze posterior samples for an architecture with 32, 64 and 128 neurons per layer in Figure \ref{fig: CT posterior samples} and Table \ref{tab: CT}. We see that UCoS and conditional method have a significant edge in terms of computational time over the other methods as a consequence of avoiding repeated forward evaluations during sampling. With larger number of neurons per layer, conditional method produces relatively similar samples as UCoS in terms bias and standard deviation, but has worse statistics and visual sample quality when using only 32 neurons. Due to the complexity of the score approximation, the performance of conditional method at 32 nodes per layer differs the most from the performance using more nodes.
Sample quality of UCoS seems comparable or better to the unconditional methods (SDE ALD, DPS and Proj) in terms of standard deviation and bias and a naive visual inspection. Noticeably, UCoS generates samples with low variance compared to the other methods.

 \begin{figure}[t]
     \centering
     % First column: Two vertically stacked images
     \begin{subfigure}[c]{0.07\textwidth}
        \textbf{32}
        \text{nodes}
     \end{subfigure}
     \begin{subfigure}[c]{0.12\textwidth}
         \centering
         truth \\
         \vspace{1mm}
         \includegraphics[scale=0.05]{Images/CT_imaging/32_nodes/ground_truth.png} \\ % First image
         measurement \\
         \vspace{1mm}
         \includegraphics[scale=0.09]{Images/CT_imaging/32_nodes/measurement.png}   % Second image
     \end{subfigure}
     % Second column: Single image
     \begin{subfigure}[c]{0.5\textwidth}
         \centering
         \hspace{1.5cm} samples \hspace{2.5cm} bias \hspace{5mm} std
         \includegraphics[width=\textwidth]{Images/CT_imaging/32_nodes/summary325.png}
     \end{subfigure}   
    \begin{subfigure}[c]{0.2\textwidth}
        \vspace{8mm}
        SDE ALD 
        \vspace{7mm}
        \\
        \vspace{7mm}
        DPS \\
        \vspace{7mm}
        Proj \\
        \vspace{8.5mm}
        Conditional \\
        \vspace{7mm}
        UCoS 
    \end{subfigure}
    
     \begin{subfigure}[c]{0.075\textwidth}
        \textbf{64}
        \text{nodes}
     \end{subfigure}
     \begin{subfigure}[c]{0.12\textwidth}
         \centering
         truth \\
         \vspace{1mm}
         \includegraphics[scale=0.05]{Images/CT_imaging/64_nodes/ground_truth.png} \\ % First image
         measurement \\
         \vspace{1mm}
         \includegraphics[scale=0.09]{Images/CT_imaging/64_nodes/measurement.png}   % Second image
     \end{subfigure}
     % Second column: Single image
     \begin{subfigure}[c]{0.5\textwidth}
         \centering
         \includegraphics[width=\textwidth]{Images/CT_imaging/64_nodes/summary645.png}
     \end{subfigure}   
\begin{subfigure}[c]{0.2\textwidth}
        \vspace{6mm}
        SDE ALD 
        \vspace{7mm}
        \\
        \vspace{7mm}
        DPS \\
        \vspace{7mm}
        Proj \\
        \vspace{7mm}
        Conditional \\
        \vspace{7mm}
        UCoS 
    \end{subfigure}
    
     \begin{subfigure}[c]{0.07\textwidth}
        \textbf{128}
        \text{nodes}
     \end{subfigure}
    \begin{subfigure}[c]{0.12\textwidth}
        \centering
        truth \\
        \vspace{1mm}
        \includegraphics[scale=0.05]{Images/CT_imaging/128_nodes/ground_truth.png} \\ % First image
        measurement \\
        \vspace{1mm}
        \includegraphics[scale=0.09]{Images/CT_imaging/128_nodes/measurement.png}   % Second image
    \end{subfigure}
    % Second column: Single image
    \begin{subfigure}[c]{0.5\textwidth}
        \centering
        \includegraphics[width=\textwidth]{Images/CT_imaging/128_nodes/summary1285.png}
    \end{subfigure}        
   \begin{subfigure}[c]{0.2\textwidth}
        \vspace{6mm}
        SDE ALD 
        \vspace{7mm}
        \\
        \vspace{7mm}
        DPS \\
        \vspace{7mm}
        Proj \\
        \vspace{7mm}
        Conditional \\
        \vspace{7mm}
        UCoS 
    \end{subfigure}
    \caption{Summary of posterior samples for CT imaging problem. Top figure uses FNO with 32 nodes per layer, middle figure uses 64 nodes per layer and bottom row 128 nodes per layer.
    In each figure, on the left column, the top image is ground truth and the bottom is measurement data (sinogram). Both are fixed for all FNO architectures. On the right column, 5 posterior samples, their bias and standard deviation. Methods used from top to bottom: SDE ALD, DPS, Proj, Conditional, and UCoS.}
    \label{fig: CT posterior samples}
\end{figure}

\begin{table}[h]
\centering
    \begin{tabular}{c  c c c c c c c c c c c}
            \hline
            method &  \multicolumn{3}{c}{\textbf{bias} for k nodes} & & \multicolumn{3}{c} {\textbf{std} for k nodes}  & & \multicolumn{3}{c} {\textbf{time} for k nodes}  \\ 
            \cline{2-4}  \cline{6-8} \cline{10-12} 
            & k=32 & k=64 & k=128 &  &k=32 & k=64 & k=128 &  &k=32 & k=64 & k=128\\
            \hline
            SDE ALD	&  0.1014 & 0.1042 & 0.1044  & & 0.0741 & 0.1064 & 0.0644 &  & 35 & 48 & 77\\
            DPS	 &  0.0605 & 0.0552 & 0.0788 & &0.068 & 0.0768 & 0.1723 & & 69 &99 & 176\\
            Proj &  0.0644 & 0.0610 & 0.0616 && 0.0667 & 0.078 & 0.0702 & & 224 & 237 & 266 \\
            Conditional	&  0.2336 & 0.0621 & \textbf{0.0467} && 0.0923 & 0.0594 & 0.0516 & & \textbf{15} & \textbf{27} & \textbf{57}\\
            UCoS &  \textbf{0.0489} & \textbf{0.0498} & 0.0541 & & \textbf{0.0379} & \textbf{0.0316} & \textbf{0.0262} & & \textbf{15} & \textbf{27} & \textbf{57}\\
            \hline
    \end{tabular}
\caption{Summary statistics of all methods for the CT imaging problem. Left table: bias, middle: std and right table: computational time in minutes. Bias and std are computed on 1000 generated posterior samples and averaged over the number of pixels. In each table, the columns correspond to the number of neurons used per layer.}
\label{tab: CT}
\end{table}

In addition to the previous experiment, for UCoS and conditional method, we train models for numerous number of neurons on a reduced $64 \times 64$ resolution. In Figure \ref{fig: graph nodes vs error}, we compare the number of nodes to the resulting posterior bias and standard deviation for the full $256 \times 256$ resolution (crosses) and the reduced $64 \times 64$ resolution (lines).
We observe that, especially in terms of standard deviation, UCoS achieves saturation with lower model complexity. To provide context, halving number of nodes approximately halves the online sampling time.


\begin{figure}[h]
    \centering
    % First subfigure
    \begin{subfigure}[b]{0.49\linewidth} % Adjust width as needed
        \centering
        \includegraphics[width=0.7\linewidth]{Images/CT_imaging/graph/graphbias.png}
    \end{subfigure}
    \hfill % Add spacing between figures
    % Second subfigure
    \begin{subfigure}[b]{0.49\linewidth} % Adjust width as needed
        \centering
        \includegraphics[width=0.7\linewidth]{Images/CT_imaging/graph/graphstd.png}
    \end{subfigure}
    \caption{Comparison of the error dependence of the conditional method and UCoS on the parametrization of the score approximation. X-axis: number of nodes per layer in the FNO, Y-axis: L2 norm of the bias (left) or std (right), averaged over the number of pixels. Dashed lines correspond to $64 \times 64$ resolution and crosses correspond to the problem with $256 \times 256$ resolution.}
    \label{fig: graph nodes vs error}
\end{figure}

\paragraph{Deblurring}
In this example, the forward operator A corresponds to a two-dimensional convolution with a Gaussian kernel over the image domain. We observe the output function on the same grid with an additive Gaussian noise vector.

The training data is given by the Large-scale CelebFaces Attributes (CelebA) Dataset \cite{liu2015faceattributes}, which contains $> 200.000$ face images of celebrities. We scale the images to square resolution $175 \times 175$.

\begin{table}[H]
    \centering
    \begin{tabular}{c  c c c}
            \hline
             method & bias& std & time\\
             \hline
             SDE ALD & 0.0878 & 0.0915 & 300 \\
             
             DPS & 0.0872 & 0.0801 &  239\\
 
             Proj & \textbf{0.0785} & 0.0622 & 15822 \\

             Conditional & 0.2178 & 0.0950 &  \textbf{85}\\
             UCoS & 0.0884 & \textbf{0.0345} & \textbf{85} \\
             \hline
        \end{tabular}
    \caption{Summary statistics of all methods for the Deblurring problem. Bias and std are computed on 1000 generated posterior samples and averaged over the number of pixels. Computational time is given in minutes.}
    \label{tab: deblurring}
\end{table}


We demonstrate in Figure \ref{fig:deblurr posterior samples 128} and Table \ref{tab: deblurring} that UCoS and conditional achieve the fastest sampling speed. For the chosen FNO parameterization, the conditional method produces samples with higher bias and standard deviation as well as low visual quality. Visually, the posterior ensemble $(N = 1000)$  produced by UCoS is consistent with the ground truth, while other methods introduce artifacts that are inconsistent with human facial features. In terms of bias, UCoS is at a roughly similar or slightly better level than unconditional methods, while the pixel-wise standard deviation is noticeably lower than other methods.

\begin{figure}[h]
    \centering
    % First column: Two vertically stacked images
    \begin{subfigure}[c]{0.2\textwidth}
    \centering
    truth
        \includegraphics[width=\textwidth]{Images/Deblurring/truth.jpg} \\ % First image
        \vspace{5pt} % Optional spacing between the two images
        measurement
        \includegraphics[width=\textwidth]{Images/Deblurring/measurement.jpg}   % Second image
    \end{subfigure}
    % Second column: Single image
    \begin{subfigure}[c]{0.6\textwidth}
        \centering
        \hspace{1.5cm} samples \hspace{3cm} bias \hspace{5mm} std
        \includegraphics[width=\textwidth]{Images/Deblurring/summary21285.png}
    \end{subfigure}
    \begin{subfigure}[c]{0.15\textwidth}
        SDE ALD \\
        \vspace{6mm}
        \\
        DPS \\
        \vspace{6mm}
        \\
        Proj \\
        \vspace{6mm}
        \\
        Conditional \\
        \vspace{2mm}
        \\
        UCoS 
    \end{subfigure}
   \caption{Summary of posterior samples for the Deblurring problem. On the left column, the top image is ground truth and the bottom is measurement data. On the right column, 5 posterior samples and their bias and standard deviation, with methods used from top to bottom: SDE ALD, DPS, Proj, Conditional, and UCoS.}
    \label{fig:deblurr posterior samples 128}
\end{figure}



\section{Conclusion, limitations and future work}
%
This paper addresses the challenge of balancing the computational cost between offline training and online posterior sampling in large-scale inverse problems. In many such problems, evaluating the forward operator during sampling constitutes a major computational bottleneck. We introduce a novel, theoretically grounded method UCoS that entirely removes this burden during sampling by shifting it to the offline training phase, without introducing any approximation error. Our approach combines the computational benefits of the conditional method introduced in \citet{baldassari2024conditional} with a more scalable score-learning formulation. Furthermore, we demonstrate that this method is rigorous in an infinite-dimensional setting and, therefore, independent of discretization.

\fabian{We validate the correctness of our approach with numerical experiments related to CT imaging and Deblurring. UCoS generates samples with approximately comparable or better quality than unconditional approaches while offering tangible computational benefits during sampling. However, it remains task-dependent, requiring re-training when the forward problem is modified. \lam{The quality of posterior samples is assessed using simple statistics such as bias, pixel-wise standard deviation, and visual inspection. Hence, the results should be viewed as proof of concept. Further experiments are needed for a more comprehensive comparison of sample quality.}
While UCoS does not accelerate posterior sampling compared to the conditional method when using score approximation with similar complexity, we show that conditional methods, unlike UCoS, involve more complex function approximations. This complexity can hinder performance, particularly for neural operators with a limited number of parameters.}

The general principle of transforming a task-dependent unconditional score into a conditional score, as developed in this work, extends beyond the specific OU diffusion process considered here. This opens up intriguing questions about selecting the appropriate process to achieve a more efficient balance of computational effort.

%

\duclam{That said, the theoretical foundation of UCoS relies critically on the interplay between a linear forward model and a Gaussian likelihood, which together enable the diffused posterior to be reinterpreted as a diffused prior under a new, task-specific process. This structure allows us to derive an exact identity between the conditional score and a transformed unconditional score, as captured in Theorem~\ref{thm:score_id1 infinite dimension}. However, extending this exact correspondence beyond settings where one can `complete the square' -- as done in \Eqref{eq:qt_identity} -- appears nontrivial. Investigating whether similar structures can be identified in nonlinear or non-Gaussian problems remains an important avenue for future work.
}

\section*{Acknowledgements}

This work has been supported by the Research Council of Finland (RCoF) through the \emph{Flagship of advanced mathematics for sensing, imaging and modelling} (FAME), decision number 358944. Moreover, TH was supported through RCoF decision numbers 353094 and 348504. M.L. was partially supported by PDE-Inverse project of the European Research Council of the European Union,  the FAME and Finnish Quantum flagships and the grant 336786 of the RCoF. Views and opinions expressed are those of the authors only and do not necessarily reflect those of the European Union or the other funding organizations. Neither the European Union nor the other funding organizations can be held responsible for them. MVdH gratefully acknowledges the support of the Department of Energy BES, under grant DE-SC0020345, Oxy, the corporate members of the Geo-Mathematical Imaging Group at Rice University and the Simons Foundation under the MATH + X Program. The authors acknowledge the National Cancer Institute and the Foundation for the National Institutes of Health, and their critical role in the creation of the free publicly available LIDC/IDRI Database used in this study.

\bibliography{references}
\bibliographystyle{tmlr} 

%%%% Appendix %%%
\newpage
\appendix

\section{Probability measures on Hilbert spaces}
\label{Appendix:prob_measures}

\subsection{Gaussian random processes in an infinite-dimensional Hilbert space} %temporary title

This section introduces notations and outlines some basic properties of probability measures on Hilbert spaces. For a more comprehensive introduction, we refer to \citet{da2014stochastic,hairer2023introductionSPDE}.

\paragraph{Gaussian measures on Hilbert space}
Let $(H,  \langle \cdot, \cdot \rangle_H)$ be a separable Hilbert space with norm $\|\cdot\|_H = \sqrt{ \langle \cdot, \cdot \rangle_H}$. 
A bounded linear operator $C: H \to H$ is called \emph{self-adjoint} if $\langle x, C y \rangle_H = \langle C x, y \rangle_H $ for all $x,y \in H$ and \emph{positive definite} if $\langle C x, x \rangle_H \ge 0 $ for all $x\in H$. We say that a self-adjoint and positive definite operator $C$ is of \emph{trace class} if  
\begin{equation*}
    \tr_H(C):= \sum_{n=1}^\infty \langle Ce_n, e_n \rangle < \infty,
\end{equation*}
where $\{e_n\}$ is an orthogonal basic of $H$. We denote by $L_1^+(H)$ the space of all self-adjoint, positive definite and trace class operators on $H$.
A random variable taking values in $H$ is called \emph{Gaussian} if the law of $\langle h, X \rangle_H $ is Gaussian for each $h\in H$. Gaussian random variables are determined by their \emph{mean} $m = \mathbb E[X] \in H$ and their \emph{covariance operator} defined as
\begin{equation*}
    \langle C g, h \rangle = \mathbb E \left[\langle g, X - m \rangle\langle h, X - m \rangle\right].
\end{equation*}
In this case, we denote $X \sim \mathcal{N}(m, C)$. If $m=0$, $X$ is called \emph{centred}. It can be shown that if $X$ is a Gaussian random variable on $H$ then $C \in L_1^+(H)$, moreover, $\mathbb E[\|X\|_H^2] = \tr_H(C)$.

\paragraph{The Cameron--Martin space} We define the Cameron--Martin space associated with a Gaussian measure $\mu = \mathcal{N}(0, C)$ on $H$ to be the intersection of all linear spaces of full measure under $\mu$, and denote it by $H_\mu$ or $H_C$. It can be shown that $H_\mu = C^{1/2} H$ and $H_\mu$ is compactly embedded and dense in $H$. In infinite dimensions, it is necessarily the case that $\mu(H_\mu) = 0$. Moreover, $H_\mu$ can be endowed with a Hilbert space structure with an inner product
\begin{equation*}
    \langle g, h\rangle_{H_\mu} = \langle g, C^{-1}h\rangle_H = \langle C^{-1/2}g, C^{-1/2}h\rangle_H.
\end{equation*}
\begin{example}
    Let $H=\R^d$ and $\mu = \mathcal{N}(0,C)$ be a Gaussian measure on $H$ with a positive definite covariance matrix $C \in \R^{d\times d}$. Then since $C^{1/2}H = H$, the Cameron--Martin space is the whole space $\R^d$. 
\end{example}
% \begin{example}
%     Let $H = L^2([0,1])$ and consider the Wiener process $B = (B_t: t\in [0,1])$ taking values in $H$ and satisfying $B_t(0)=0$. In non-rigorous manner, the probability density function of the Wiener process is formally written as $\pi_B(b)=c\exp(-\|\partial_tb\|_{L^2(0,1)}^2/2)$. The Cameron--Martin space of $B$ is given as  $\{b(t)\in H^1([0,1]): b(0)=0\}$, where  $H^1([0,1])$ is the Sobolev space of weakly differentiable functions on $[0,1]$. It is well-known that sample paths of $B$ are almost surely nowhere differentiable, thus they do not lie in $H^1([0,1])$ almost surely. Therefore $\mu(H^1([0,1])) = 0$ where $\mu$ is the Wiener measure that describes $B$.
% \end{example}
\paragraph{Cameron--Martin's theorem} The Cameron–Martin space $H_\mu$ plays a special role in that it characterises precisely those directions in which one can translate the Gaussian measure $\mu$ without changing its null sets, thus $\mu_h:= \mathcal{N}(h, C)$ and $\mu = \mathcal{N}(0, C)$ are equivalent if and only if $h\in H_\mu$. Moreover, the Radon--Nikodym derivative of $\mu_h$ with respect to $\mu$ is given by
\begin{equation}\label{CM-theorem}
    \frac{\dd\mu_h}{\dd\mu}(x) = \exp\left( \langle h, x\rangle_{H_\mu} - \frac{1}{2}\|h\|_{H_\mu}^2 \right), \quad \mu\text{-a.s.}, x\in H.
\end{equation}
Note that since $H_\mu=C^{1/2}H$ is dense in $H$, the random variable $\langle h, x\rangle_{H_\mu} = \langle C^{-1/2}h, C^{-1/2}x\rangle_H$, $x\in H$, can be defined properly using a limiting process, see Remark 2.24 in  \citet{da2014stochastic}.





\subsection{Wiener processes $B_t$ in an infinite-dimensional Hilbert space}
\label{subsec: Wiener process definition}
\label{appendix:wiener_process}

In this subsection, we recall the definition of the Wiener processes $B_t$ in an infinite-dimensional Hilbert space.  As $H$ is an infinite-dimensional Hilbert space, the process 
$B_t$ can not be defined as random variable taking values in the Hilbert space $H$, but we need
to start our considerations with a larger Hilbert space $Y$ having the additional properties
that are described below.
At each time $t\in [0,T]$ the process $B_t$ is defined to be a Gaussian random
variable taking values in the Hilbert space $Y$ so that
 the expectation of $B_t$ is zero and the covariance operator  is
$tC_Y$, where $C_Y\colon Y\to Y$ is a symmetric and injective trace-class operator in $Y$.
That is, for $\phi,\psi\in Y$ and $0\le t<t+s\le T$ we have
\begin{eqnarray*}
& &    \mathbb E(\langle B_t, \phi\rangle_Y)=0,\\
& &\mathbb E( \langle B_t, \phi\rangle_Y\cdot 
\langle B_t, \psi\rangle_Y)=t\langle \phi,C_Y \phi\rangle_Y,
\\
& &\mathbb E( \langle B_{t+s}-B_t, \phi\rangle_Y\cdot\langle B_{t+s}-B_t, \psi\rangle_Y)=s\langle \phi,C_Y \phi\rangle_Y.
\end{eqnarray*}

We require that $Y$ and $C_Y$ are such that
\begin{equation*}
H=\hbox{Ran}(C_Y^{1/2})=C_Y^{1/2}(Y).
\end{equation*}
Under the above assumptions, $H$ is the Cameron-Martin 
of the Gaussian random variable $B_t$ for $t>0$. Then, we can define for $b\in Y$ and $\phi\in C_Y(Y)$ an
extension of $H$-inner product $\langle u, \phi\rangle_H$
by setting
$$
\langle u, \phi\rangle_H:=\langle u,C_Y^{-1} \phi\rangle_Y.
$$
Using these definitions,
it holds that for all $\phi,\psi\in C_Y(Y)\subset H$ and $0\le t<t+s\le T$ we have
\begin{eqnarray*}
& &    \mathbb E(\langle B_t, \phi\rangle_H)=0,\\
& &\mathbb E( \langle B_t, \phi\rangle_H\cdot 
\langle B_t, \psi\rangle_H)=t\langle \phi, \phi\rangle_H,
\\
& &\mathbb E( \langle B_{t+s}-B_t, \phi\rangle_H\cdot 
\langle B_{t+s}-B_t, \psi\rangle_H)=s\langle \phi, \phi\rangle_Y.
\end{eqnarray*}
Motivated by these formulas, we call $B_t$, $0\le t\le T$, a Wiener process in $H$
having the (generalized) covariance operator $tI$.

We also note that in formula \Eqref{eq: diffusion model; initiated from posterior}, we use the symmetric trace-class operator $\mathcal{C}\colon H \to H$ and the increments $\mathcal{C}^{1/2} dB_t$. These increments can be interpreted as the differences of the random process $t \mapsto \mathcal{C}^{1/2} B_t$. At any time $0 \le t \le T$, the random variable $\mathcal{C}^{1/2} B_t$ takes values in the space $H$, and its covariance operator in $H$ is $\mathcal{C}\colon H \to H$.


\section{Proofs for Section \ref{sec:UCoS}}

% \subsection{Completing the square}
% \label{sec:completing_square}
% \begin{proof}[Proof of \Eqref{eq:qt_identity}] 
% We first write
% \begin{equation}
%     \label{eq:Az-y}
%     \begin{aligned}
%     \norm{Az-y}^2_\Gamma = & \left\langle \Gamma^{-1}(Az - y), Az - y \right \rangle \\
%     = & z^\top A^\top \Gamma^{-1} Az - z^\top A^\top \Gamma^{-1} y - y^\top \Gamma^{-1} Az + y^\top \Gamma^{-1} y.
%     \end{aligned}
% \end{equation}
% Now since the transitional probability for the multivariate OU process is given by $\rho_t(x | x_0) = {\mathcal N}(e^{-t/2}x_0, (1-e^{-t})\C)$, we have
% \begin{equation}
%     \label{eq:transition_prop}
%     \begin{aligned}
%         \rho_t(z | x) = & \frac{1}{(2\pi)^{n/2}[\det \C (1 - e^{-t})]^{1/2}} \exp \left( -\frac{1}{2} \left(z - xe^{-t/2} \right)^\top \C^{-1}(1 - e^{-t})^{-1}\left(z - xe^{-t/2} \right) \right) \\
%         = & \frac{1}{(2\pi)^{n/2}[\det \C (1 - e^{-t})]^{1/2}} \exp \left( -\frac{1}{2(1-e^{-t})} \left(z^\top e^{-t}\C^{-1} z - z^\top e^{-t/2}\C^{-1} x - x^\top \C^{-1} z e^{-t/2} + x^\top \C^{-1} x \right) \right)
%     \end{aligned} 
% \end{equation}
% By combining \Eqref{eq:Az-y} and \Eqref{eq:transition_prop},
% \begin{equation*}
%     \begin{aligned}
%         \exp\left(-\frac 12 \norm{Az-y}^2_\Gamma\right) \rho_t(x | z) \propto \exp\left(-\frac{1}{2}\left( z^\top \Sigma_t^{-1} z - z^\top \left(A\top \Gamma^{-1}y + \frac{1}{e^{t/2} - e^{-t/2}} \C^{-1}x \right) \right) - \left(y^\top \Gamma^{-1}A + \frac{1}{e^{t/2} - e^{-t/2}} x^\top \C^{-1} \right) \right)
%     \end{aligned}
% \end{equation*}
% TBC
% \end{proof}

\subsection{Task-depedent score} %temporary title

Recall that $C_t = \Gamma + (e^t -1 ) A \C A^* : H\to H$ and $\Sigma_t = (e^t -1)\C - (e^t-1)^2\C A^* C_t^{-1} A \C : H\to H$.

\begin{proof}[Proof of Lemma \ref{lem: Sigma_t is covariance operator}]
It is clear that $\Sigma_t$ is self-adjoint. For positive definiteness, we write
    \begin{eqnarray*}
        \Sigma_t 
        %& = & (e^t -1)\C - (e^t-1)^2\C A^* C_t^{-1} A \C \\
        &=& \C^{1/2} \left((e^t-1)I - (e^t-1)^2  \C^{1/2} A^* C_t^{-1}A \C^{1/2}\right)  \C^{1/2} \\
        & = & \C^{1/2} \left( \frac{1}{e^t-1}I + \C^{1/2}A^* \Gamma^{-1}A \C^{1/2}\right)^{-1} \C^{1/2},
    \end{eqnarray*}
where we applied Lemma \ref{lem: Useful equality that show that Sigma is still p.d.} (a) for the last identity.
Since $\C A^* C_t^{-1} A \C$ is also positive definite, it immediately follows that
\begin{equation*}
    0<\tr_H(\Sigma_t) = (e^t-1) \tr_H(\C) - (e^t-1)^2 \tr_H(\C A^* C_t^{-1} A \C \leq (e^t-1) \tr_H(\C) < \infty
\end{equation*}
and, consequently, $\Sigma_t$ is trace-class.
\end{proof}

\begin{proof}[Proof of Lemma \ref{lem: Score function of posterior process in infinite dimension}] 
We replicate the proof of Lemma 1 in \citet{pidstrigach2023infinitedimensional} for the process $\widetilde{X}^{\mu}_t$.
The density $\tilde p_t(\cdot |x_0)$ of $\widetilde{X}^{\mu}_t $ conditioned on $\widetilde{X}^{\mu}_0 = x_0$ is Gaussian centered at $x_0$ with covariance $\Sigma_t$. 

For what follows, let $\tilde{p}_0(x_0)$ be the density of $\widetilde{X}_0^{\mu}$ and $\tilde p_t(x)$ the density of $\widetilde{X}^{\mu}_t$.
We apply Leibniz's rule to obtain
\begin{eqnarray*}
    \Sigma_t \nabla_x \log \tilde p_t(x) 
%    & = & \Sigma_t  \frac{\nabla_x \tilde p_t(x)}{\tilde p_t(x)}\\
    &=&  \Sigma_t \left(\frac{1}{\tilde p_t(x)} \nabla_x \int \tilde p_t(x|x_0) d \tilde p_0(x_0)\right) \\
    & = &-\Sigma_t  \int \Sigma_t^{-1}(x-x_0) \frac{\tilde p_t(x| x_0)}{\tilde p_t(x)} d \tilde p_0(x_0) \\
    & = & - \int (x-x_0) \tilde p(x_0 |x) \\
    & = &
    - \left(x - \E(\widetilde{X}^{\mu}_0 |\widetilde{X}^{\mu}_t = x)\right), 
\end{eqnarray*}
where we utilized Bayes' formula for the second last identity.
This proves the claim.
\end{proof}



\begin{lemma}\label{lem: Useful equality that show that Sigma is still p.d.}
For any $t>0$, we have that
\begin{itemize}
    % \item [a)] 
    % \begin{eqnarray*}
    %     \left( I -(e^t-1) A^*C_t^{-1}A \C \right)^{-1} = I + (e^t-1)A^* \Gamma^{-1}A \C
    % \end{eqnarray*}
    \item[(a)] the linear operator $\Xi_t = (e^t-1)I -(e^t-1)^2 \C^{1/2} A^*C_t^{-1}A \C^{1/2} : H \to H$ is bijective and
    \begin{eqnarray*}
    \Xi_t^{-1} = \frac{1}{e^t-1}I + \C^{1/2} A^* \Gamma^{-1}A \C^{1/2},
    \end{eqnarray*}
    \item[(b)] the linear operator $\Xi_t' = I -(e^t-1) \C A^*C_t^{-1}A$ is bijective and
    \begin{eqnarray*}
        \left(\Xi_t'\right)^{-1} = I + (e^t-1)\C A^* \Gamma^{-1}A; \quad \text{and}
    \end{eqnarray*}
    \item[(c)] it holds that
    \begin{eqnarray*}
        (e^t-1) \C A^* C_t^{-1} = \Sigma_t A^* \Gamma^{-1}.
    \end{eqnarray*}
\end{itemize}
\end{lemma}

\begin{proof}
a) Invertibility from the right can be derived by straightforward computation
    \begin{multline*}
         \Xi_t \bigg( \frac{1}{e^t-1}I + \C^{1/2}A^* \Gamma^{-1}A \C^{1/2}\bigg) \\
%        & = I + (e^t-1) \C^{1/2}A^* \Gamma^{-1}A \C^{1/2} \\
%        &\quad - (e^t-1) \C^{1/2} A^* C_t^{-1}A \C^{1/2} - (e^t-1)^2  \C^{1/2} A^* C_t^{-1}A \C A^* \Gamma^{-1}A \C^{1/2} \\
         = I + (e^t-1) \C^{1/2}A^* \Gamma^{-1}A \C^{1/2} 
         -  (e^t-1)\C^{1/2} A^* C_t^{-1} \underbrace{\left((e^t-1) A\C A^* + \Gamma \right)}_{= C_t} \Gamma^{-1}A\C^{1/2} = I
    \end{multline*}
Similarly, invertibility from the left follows from an analogous computation. The invertibility in (b) can be established using the same arguments.
    % In the same fashion:
    % \begin{eqnarray*}
    %     & & \left( \frac{1}{e^t-1}I + \C^{1/2}A^* \Gamma^{-1}A \C^{1/2}\right) \left((e^t-1)I - (e^t-1)^2  \C^{1/2} A^\top C_t^{-1}A \C^{1/2}\right)  \\ 
    %     & = & I  - (e^t-1) \C^{1/2} A^* C_t^{-1}A\C^{1/2} + (e^t-1) \C^{1/2} A^* C_t^{-1}A\C^{1/2} \\
    %     & & - (e^t-1)^2 \C^{1/2}A^* \Gamma^{-1}A\C A^* C_t^{-1}A \C^{1/2} \\
    %     & = & I + (e^t-1) \C^{1/2} A^* C_t^{-1}A\C^{1/2} \\
    %     & & - (e^t-1) \C^{1/2}A^* \Gamma^{-1} \underbrace{\left( \Gamma + (e^t-1) A^*\C A) \right)}_{= C_t}C_t^{-1} A\C^{1/2} \\
    %     &=& I
    %     \end{eqnarray*}
For the identity in (c), we have
    \begin{eqnarray*}
        \Sigma_t A^* \Gamma^{-1} C_t 
%        & = & \Sigma_t A^* \Gamma^{-1}(\Gamma + (e^t-1)A\C A^* ) \\
        & = & (e^t-1)\C A^* -(e^t-1)^2 \C A^*C_t^{-1} A \C A^* + (e^t-1)^2 \C A^* \Gamma^{-1} A \C A^* \\
        & & - (e^t-1)^3 \C A^* C_t^{-1} A \C A^* \Gamma^{-1} A \C A^* \\
        & = & (e^t-1) \C A^* + (e^t-1)^2 \C A^* \Gamma^{-1} A \C A^* \\
        & & - (e^t-1)^2 \C A^* C_t^{-1} \underbrace{\left((e^t-1)A \C A^* + \Gamma \right)}_{=C_t} \Gamma^{-1} A \C A^* \\
        & =& (e^t-1) \C A^*.
    \end{eqnarray*}
The desired identity follows by inverting $C_t$.
\end{proof}

%Before we start proving the main Theorem \ref{thm:score_id1 infinite dimension}, we recall some necessary elements of theory of probability on general Hilbert spaces.


\subsection{Proof of Theorem \ref{thm:score_id1 infinite dimension}} 
We follow the structure in Section \ref{sec: finite dimensional UCoS} and first prove the following equivalent of \Eqref{eq:finite_dim_score_id 2}. Theorem \ref{thm:score_id1 infinite dimension} follows directly from substituting  \Eqref{eq: r equation} into Theorem \ref{thm:prestep to score_id1 infinite dimension} while utilizing Lemma \ref{lem: writing m_t for no forward}.
\begin{theorem}\label{thm:prestep to score_id1 infinite dimension}
Let $H_{\C}$ be the Cameron-Martin space of $\C$ and assume that the prior satisfies $\mu(H_{\C}) = 1$. For $t>0$, let $\Sigma_t$ be given by  \Eqref{eq: formal definition of Sigma_t} and define   
\begin{eqnarray}\label{eq: definition of mut}
    m_t(x, y) = e^{t/2}x + (e^t - 1)\C A^* C_t^{-1}(y - e^{t/2}Ax).
\end{eqnarray}
Then the conditional score function $s(x,t;\mu^y)$ related to it holds that
\begin{equation}
    \label{eq:conditional_score_id}
    s(x, t; \mu^y) = \lambda(t) \left(  \Tilde{s}\left(m_t(x, y\right), t; \mu) + m_t(x, y) - e^{t/2}x \right)
\end{equation}
for $(x,y) \in H \times \R^m$ a.e. in $\L(X_t^\mu, Y)$ and $t>0$.
\end{theorem}


Theorem \ref{thm:prestep to score_id1 infinite dimension} is a direct consequence of the following proposition.
\begin{proposition}\label{prop:identity_of_expectations}
    Let the assumptions of Theorem \ref{thm:score_id1 infinite dimension} hold. Then
    \begin{eqnarray}\label{eq: equality of expectations in infinite dimensions}
       \E(\widetilde{X}^{\mu}_0 |\widetilde{X}^{\mu}_t = m_t(x,y)) =  \E(X^{\mu}_0 |Y = y, X^{\mu}_t = x),
    \end{eqnarray}
    for $(x,y) \in H \times \R^m$ a.e. in $\L(X_t^\mu, Y)$ and $t>0$.
\end{proposition}

%    \begin{eqnarray*}
%            -\frac{1}{1-e^{-t}}\left(x - e^{-t/2}\E(X^{\mu^y}_0 |X^{\mu^y}_t = x)\right) &=& \frac {1}{(e^{t/2} - e^{-t/2})}\left( \left(-m_t + \E(\widetilde{X}^{\mu}_0 |\tilde{X^{\mu}}_t = m_t)\right) + m_t - e^{t/2}x \right).
%    \end{eqnarray*}
%    For the terms involving $m_t$ and $x$ notice that $\frac{e^{t/2}}{e^{t/2} - e^{-t/2}} = \frac{1}{1-e^{-t}}$ to see that they equal.
%\end{proof}

Throughout this section, we denote the transition kernel densities
\begin{eqnarray*}
        n_t(x_0, x) := \frac{d \N( x_0, (e^t-1) \C)}{d \N(0, (e^t-1) \C)}(e^{t/2}x), \quad 
        \tilde n_t(x_0, x) := \frac{d \N(x_0, \Sigma_t)}{d \N(0, \Sigma_t)}(x),
    \end{eqnarray*}
whenever the Radon--Nykodym derivatives make sense.

A plan of proof for Proposition \ref{prop:identity_of_expectations} is as follows. We first develop auxiliary results: Lemma \ref{lem: training loss} shows that the laws of $m_t(X^\mu_t, Y)$ and $\widetilde{X}^\mu_t$ coincide when conditioned on $X_0=0$. Lemma \ref{lem: tranisition kernels} is used to express each expectation in \Eqref{eq: equality of expectations in infinite dimensions} in terms of transition kernels $n_t$ and $\tilde n_t$ of the corresponding forward SDEs, which in turn can be written as Radon--Nykodym derivatives of certain measures. After that, we use Lemma \ref{lem: relation between CM spaces of U and Vt} to show that the measures $\N(0, (e^t-1)\C)$ and $\N(0, \Sigma_t)$ are equivalent, in particular their Cameron--Martin spaces equal, which concludes the proof.
Finally, we put together the argument at the end of the section.

\begin{lemma}\label{lem: training loss}
Let $Z_1 \sim \N(0, (1-e^{-t})\C)$, $Z_2 \sim \N(0,\Gamma)$ and $Z_3 = \N(0,\Sigma_t)$ be mutually independent Gaussian random variables on $H$.
Moreover, let $x_0 \in H$ be arbitrary. It holds that
\begin{equation*}
	\L(m_t(e^{-t/2}x_0 + Z_1, Ax_0 + Z_2)) = \L(x_0 + Z_3).
\end{equation*}
for any $t>0$.
\end{lemma}
\begin{proof}
	We have that
    \begin{eqnarray*}
         m_t(e^{-t/2}x_0 + Z_1, Ax_0 + Z_2) & = & x_0 + e^{t/2}Z_1 + (e^t-1) \C A^* C_t^{-1}(Z_1 - e^{t/2}A Z_1) \\
         & = & x_0 + e^{t/2} \frac{1}{e^t-1}\Sigma_t \C^{-1} Z_1 + \Sigma_t A^* \Gamma^{-1} Z_2
    \end{eqnarray*}
    is a Gaussian random variable centered at $x_0$ with a covariance
    \begin{eqnarray*}
        \text{Cov}(m_t(e^{-t/2}x_0 + Z_1, Ax_0 + Z_2)) &=& \Sigma_t A^* \Gamma^{-1} A \Sigma_t + \frac{1}{e^t-1} \Sigma_t \C^{-1} \Sigma_t \\
        & = & \Sigma_t \left( A^* \Gamma^{-1}A + \frac{1}{e^t-1} \C^{-1}\right) \Sigma_t = \Sigma_t.
    \end{eqnarray*}
    This proves the claim.
\end{proof}

%The following lemmas are used in the proof of Proposition \ref{prop:identity_of_expectations}.
\begin{lemma}\label{lem: tranisition kernels}
    The following holds
    \begin{itemize}
    \item[(i)] For $(x,y) \in H \times \R^m$ a.e. in $\L(X_t^\mu, Y)$, 
    \begin{eqnarray*}
        \E(X^{\mu}_0 | Y = y, X^{\mu}_t = x) = \frac{\int_H x_0 n_t(x_0, x) \mu^y(dx_0)}{\int_H n_t(x_0, x) \mu^y(dx_0)}.
    \end{eqnarray*} 
    % with 
    % \begin{eqnarray*}
    %     n_t(x_0, x) 
    %     &=& \exp\left( e^{t/2} \langle x_0, x \rangle_{H_{\widetilde{\C}_t}} - \frac{1}{2}\norm{x_0}_{H_{\widetilde{\C}_t}}^2 \right).
    % \end{eqnarray*}
    \item[(ii)] For $x\in H$ $\L(\widetilde{X}_t^{\mu})$-a.e.,
    \begin{eqnarray*}
        \E(\widetilde{X}^{\mu}_0 |\widetilde{X}^{\mu}_t = x) = \frac{\int_H x_0 \tilde n_t(x_0, x) \mu(dx_0)}{\int_H \tilde n_t(x_0, x) \mu(dx_0)}.
    \end{eqnarray*}
    % with 
    % \begin{eqnarray*}
    %     \tilde n_t(x_0, x) = \frac{d \N(x_0, \Sigma_t)}{d \N(0, \Sigma_t)}(x) = \exp\left( \langle x_0, x \rangle_{H_{\Sigma_t}} - \frac{1}{2}\norm{x_0}^2_{H_{\Sigma_t}} \right)
    % \end{eqnarray*}    
    \end{itemize}
\end{lemma}

\begin{proof}
{\bf (i):} We first observe that, if $X \sim \N(x_0, \C)$ then $\alpha X \sim \N(\alpha x_0, \alpha^2 S)$ for any $\alpha > 0$, and as a direct consequence
\begin{equation}
    \label{eq:multiplying_gaussian_rv}
    \N(x_0, \C)(A) = \N(\alpha x_0, \alpha^2 \C)(\alpha A),
\end{equation}
for any $\alpha > 0$ and $A\in \mathcal{B}(H)$.
By using this property for $\alpha = e^{-t/2}$, it holds that
\begin{eqnarray*}
    n_t(x_0, x) & = & \frac{d \N( x_0, (e^t-1) \C)}{d \N(0, (e^t-1) \C)}(e^{t/2}x) \\ 
    & = & \frac{d \N(e^{-t/2} x_0, (1-e^{-t}) \C)}{d \N(0, (1-e^{-t}) \C)}(x) \\
       & =& \E(X^{\mu}_0 | Y = y, X^{\mu}_t = x),
\end{eqnarray*}
for $(x,y) \in H \in \R^m$ a.e. in $\L(X_t^\mu, Y)$, where the last equality follows from the proof of Theorem 2 in %\cite[Proof of Thm. 2]
\citet{pidstrigach2023infinitedimensional} and that $\N(e^{-t/2} x_0, (1-e^{-t}) \C)$ is the transition kernel of the forward SDE in \Eqref{eq: BSDE infdim}.

{\bf (ii):} We repeat the aforementioned argument from \citet{pidstrigach2023infinitedimensional} adapted to our case.
The joint distribution of $\widetilde{X}^{\mu}_t, \widetilde{X}^{\mu}_0$ is given by
        $\tilde{n}(x_0, x)(\N(0, \Sigma_t))(dx) \otimes \mu(dx_0)$.
    Indeed, for any ${\mathcal A} \in \sigma(\widetilde{X}^{\mu}_0), {\mathcal B}\in \sigma(\widetilde{X}^{\mu}_t)$,
    \begin{eqnarray*}
        \iint_{{\mathcal A}\times {\mathcal B}} \tilde n(x_0, x) \N(0, \Sigma_t)(dx) \mu(dx_0)
        & = &  \iint_{{\mathcal A}\times {\mathcal B}}  \frac{d \N(x_0, \Sigma_t)}{d \N(0, \Sigma_t)}(x) \N(0, \Sigma_t)(dx) \mu(dx_0) \\
        & = &  \iint_{{\mathcal A}\times {\mathcal B}}  \N(x_0, \Sigma_t)(dx) \mu (dx_0) \\
        & = & \P( \widetilde{X}^{\mu}_0 \in {\mathcal A}, \widetilde{X}^{\mu}_t \in {\mathcal B}),
    \end{eqnarray*}
    where it is used that $\N(x_0, \Sigma_t)$ is the forward transition kernel of the process $\widetilde{X}^{\mu}$.
    We show that
    \begin{eqnarray*}
        f(x) = \frac{\int x_0 \tilde n_t(x_0, x) \mu(dx_0)}{\int \tilde n_t(x_0, x) \mu(dx_0)}
    \end{eqnarray*}
    is a version of the conditional expectation $\E(\widetilde{X}^{\mu}_0 |\widetilde{X}^{\mu}_t = x)$. 
    Let $\P_t$ be the law of $\widetilde{X}^{\mu}_t$, that is, 
    \begin{eqnarray*}
        \P_t(A) = \P(\widetilde{X}_t^{\mu} \in A), \quad A\in \sigma(\widetilde{X}^{\mu}_t).
    \end{eqnarray*}
    The function $f(x)$ is $\sigma(\widetilde{X}^{\mu}_t)$-measurable by Fubini's theorem and for any ${\mathcal A}\in \sigma(\widetilde{X}^{\mu}_t)$,
    \begin{eqnarray*}
        \E_{x \sim \L(X^{\mu}_t)} \left(1_{\mathcal A} f(x)\right) &=& \int_{\mathcal A} \frac{\int_H x_0 \tilde n_t(x_0, x) \mu(dx_0)}{\int_H \tilde n_t(x_0, x) \mu(dx_0)} d \P_t(x) \\
        & = & \iint_{H\times {\mathcal A}} \frac{\int_H x_0 \tilde n_t(x_0, x) \mu(dx_0)}{\int_H \tilde n_t(x_0, x) \mu(dx_0)} \tilde n(\tilde x_0, x_t) \N(0, \Sigma_t)(dx) \mu(d\tilde x_0)\\
        & = & \iint_{{\mathcal A}\times H} x_0 \tilde n_t(x_0, x) \mu(dx_0) \frac{\int_H \tilde n(\tilde x_0, x)  \mu(d\tilde x_0)}{\int_H \tilde n_t(x_0, x) \mu(dx_0)} \N(0, \Sigma_t)(dx)
        \\ & = & \iint_{{\mathcal A}\times H} x_0 \tilde n_t(x_0, x) \mu(dx_0)  \N(0, \Sigma_t)(dx) = \E_{x_0 \sim \widetilde{X}^{\mu}_0}(1_{\mathcal A} x_0).
    \end{eqnarray*}
    The above properties define the conditional expectation and we can conclude the proof.
\end{proof}
%\begin{lemma}\label{lem: funny way to write the CM space}
%    Let $S \in L_1^+(H)$ and $U$ its Cameron-Martin space. 
%    For a set A define the closure of A w.r.t a norm $\norm{\cdot}$ as $\cl^{\norm{\cdot}}(A)$.
%    We can write 
%    \begin{eqnarray*}
%        U = \cl^{\norm{\cdot}_{U}}(S(H)).
%    \end{eqnarray*}
%\end{lemma}
%\begin{proof}
%   The claim is equivalent to showing
%   \begin{eqnarray*}
%       U = \cl^{\norm{\cdot}_{U}}(S^{1/2}(U)).
%   \end{eqnarray*}
%Because $U$ is closed, $\cl^{\norm{\cdot}_U}(U) = U$ and since $S^{1/2}(U) \subset U$, 
%\begin{eqnarray*}
%    \cl^{\norm{\cdot}_U}(S^{1/2}(U)) \subset \cl^{\norm{\cdot}_U}(U) = U.
%\end{eqnarray*}
%For the other inclusion, let $x \in U$. Because $S^{1/2}(U)$ is dense in $U$, there exists some $(s_n)_{n \in \Natural} \subset S^{1/2}(U)$ such that
%\begin{eqnarray*}
%    \norm{s_n - x}_U \rightarrow 0.
%\end{eqnarray*}
%This proves that $x \in \cl^{\norm{\cdot}_U}(S^{1/2}(U))$.
%\end{proof}

\begin{lemma}\label{lem: relation between CM spaces of U and Vt}
    We have that
    \begin{itemize}
        \item [(i)] $\Sigma_t(H) = \C(H)$; moreover,
    \begin{eqnarray}\label{eq: Sigma_t_inv in infinite dimension}
        \Sigma_t^{-1}\big|_{\Sigma_t(H)} = (A^*\Gamma^{-1}A  +(e^t-1)^{-1} \C^{-1})\big|_{\C(H)}.
    \end{eqnarray}
        \item [(ii)] $H_{(e^t-1) \C} =H_{\Sigma_t}$. 
        \item [(iii)] 
        The measures $\N(0, (e^t-1)\C)$ and $\N(0, \Sigma_t)$ are equivalent.
        \item [(iv)] 
        For $x_0 \in H_{(e^t-1) \C}$, we have
        \begin{equation}\label{eq: Equation for new Cameron-Martin space}
   \langle x_0, \cdot \rangle_{H_{\Sigma_t}} = \langle x_0, A^*\Gamma^{-1}A \cdot \rangle_H +  \langle x_0, \cdot \rangle_{H_{(e^t-1) \C}}
    \end{equation}
    in $\N(0,\Sigma_t)$-a.e. 
    \end{itemize}
\end{lemma}

\begin{proof} 
For the purpose of this proof, we abbreviate
\begin{equation}\label{eq:abbrev_(e^t-1)C}
    \widetilde{\C}_t = (e^t-1) \C.
\end{equation}

{\bf (i):} We first show that $ \widetilde{\C}_t(H) = \Sigma_t(H)$.     Note that we can write $\Sigma_t$ as
    \begin{eqnarray*}
        \Sigma_t =  \widetilde{\C}_t\left( I -\widetilde{\C}_t A^*C_t^{-1}A \right),
    \end{eqnarray*}
    where the last factor on the right-hand side is invertible (see Lemma \ref{lem: Useful equality that show that Sigma is still p.d.} (b)) proving that ranges of $\Sigma_t$ and $\widetilde{\C}_t$ coincide. 
    To prove \Eqref{eq: Sigma_t_inv in infinite dimension}, let $z \in \Sigma_t(H) = \C(H)$, then for some $x\in H$, 
    \begin{eqnarray*}
        z = \Sigma_t x =  \widetilde{\C}_t x - \widetilde{\C}_t A^* C_t^{-1}A \widetilde{\C}_t x.
    \end{eqnarray*}
    We have
    \begin{eqnarray*}
        \left(A^* \Gamma^{-1} A  + \widetilde{\C}_t^{-1} \right)z & =&  A^* \Gamma^{-1} A \widetilde{\C}_t x -  A^* \Gamma^{-1} A \widetilde{\C}_t A^* C_t^{-1}A \widetilde{\C}_t x + x  - A^* C_t^{-1} A \widetilde{\C}_t x \\
        & =& x + A^* \Gamma^{-1}A \widetilde{\C}_t x - A^*\Gamma^{-1} \left(\Gamma + A \widetilde{\C}_t A^* \right) C_t^{-1}A\widetilde{\C}_t x = x,
    \end{eqnarray*}
    since $\Gamma + A \widetilde{\C}_t A^* = C_t$. This shows that
    \begin{eqnarray*}
        \Sigma_t^{-1}z = x = \left( A^*\Gamma^{-1}A  + \widetilde{\C}_t^{-1} \right)z.
    \end{eqnarray*}
{\bf (ii):} Notice first that we can write
    \begin{eqnarray*}\label{eq: funny way to write the CM space}
        H_{\widetilde{\C}_t} = \cl^{\norm{\cdot}_{H_{\widetilde{\C}_t}}}(\widetilde{\C}_t(H)), \quad H_{\Sigma_t} = \cl^{\norm{\cdot}_{H_{\Sigma_t}}}(\Sigma_t(H)),
    \end{eqnarray*}
    where $\cl^{\norm{\cdot}_{H_1}}(H_2)$ denotes the closure of $H_2\subset H_1$ w.r.t $\norm{\cdot}_{H_1}$. 
    Let us prove that the norms $\norm{\cdot}_{H_{\widetilde{\C}_t}} $ and $ \norm{\cdot}_{H_{\Sigma_t}}$ are equivalent on $\Sigma_t(H) = \widetilde{\C}_t(H)$ as, together with (i), this will prove the statement.
    
    To this end, let $x \in \Sigma_t(H) = \widetilde{\C}_t(H)$ and apply $(i)$ and Cauchy--Schwarz inequality to obtain
    %It follows from this and Cauchy--Schwarz inequality that
    \begin{eqnarray*} \norm{x}^2_{H_{\widetilde{\C}_t}} \le \norm{x}^2_{H_{\Sigma_t}}
        & = & \langle x, A^* \Gamma^{-1}A x \rangle_H  + \norm{x}_{H_{\widetilde{\C}_t}}^2 \\
        & \le &   \opnorm{\widetilde{\C}_t^{1/2}A^* \Gamma^{-1}A\widetilde{\C}_t^{1/2}}\norm{x}_{H_{\widetilde{\C}_t}}^2 + \norm{x}_{H_{\widetilde{\C}_t}}^2\\
        & =& \left( 1+  \opnorm{\widetilde{\C}_t^{1/2}A^* \Gamma^{-1}A\widetilde{\C}_t^{1/2}} \right) \norm{x}_{H_{\widetilde{\C}_t}}^2,
    \end{eqnarray*}
    where the operator norm $\opnorm{\widetilde{\C}_t^{1/2}A^* \Gamma^{-1}A\widetilde{\C}_t^{1/2}}$ is finite. 
 
    {\bf (iii):} From $(ii)$, the covariance operators of $\mu$ and $\nu$ have the same Cameron--Martin space, by the Feldman--Hajek theorem we need to show that
    $\widetilde{\C}_t^{-1/2} \Sigma_t \widetilde{\C}_t^{-1/2} - I$
    is Hilbert-Schmidt. To this end, note that
    \begin{equation*}
        \widetilde{\C}_t^{-1/2} \Sigma_t \widetilde{\C}_t^{-1/2} - I = - \widetilde{\C}_t^{1/2} A^* C_t^{-1} A \widetilde{\C}_t^{1/2} := B.
    \end{equation*}
    The operator $B^2$ is of trace class, since
    \begin{eqnarray*}
        \tr_H (B^2)&=& \tr_H\left( \widetilde{\C}_t^{1/2}A^* C_t^{-1} A \widetilde{\C}_t A^* C_t^{-1} A \widetilde{\C}_t^{1/2} \right)  \\
        & \leq & \opnorm{A^* C_t^{-1} A \widetilde{\C}_t A^* C_t^{-1} A} \tr_H(\widetilde{\C}_t) < \infty,
    \end{eqnarray*}
    proving the claim.
    
    {\bf (iv):} Let $x_0 \in \Sigma_t(H)$. By $(i)$, we have
    \begin{equation}\label{eq:inner_product_in-2CM}
    \begin{aligned}
        \langle x_0, \cdot \rangle_{H_{\Sigma_t}} & = \langle   \Sigma_t^{-1} x_0, \cdot \rangle_{H} \\
        &= \langle  A^*\Gamma^{-1}A   x_0, \cdot \rangle_H +  \langle   \widetilde{\C}_t^{-1} x_0, \cdot \rangle_H \\
        & =  \langle A^*\Gamma^{-1}A   x_0, \cdot\rangle_H + \langle x_0, \cdot \rangle_{H_{\widetilde{\C}_t}} \quad \N(0,\Sigma_t)\text{-a.e.}
    \end{aligned}
    \end{equation}
    Now since $\Sigma_t(H) = \C(H)$ is dense in $H_{\Sigma_t} = H_{\widetilde{\C}_t}$, the identity above can be uniquely extended to $H_{\Sigma_t}$ by the white noise mapping, see \citep[p. 23]{da2006introduction}. This completes the proof.
\end{proof}


\begin{corollary}
\label{cor:ID1}
We have
    \begin{equation}\label{eq:ID1}
        e^{t/2} \langle x_0, X_t^\mu \rangle_{H_{\widetilde{\C}_t}} + \langle x_0, A^* \Gamma^{-1}Y \rangle_H = \langle x_0, m_t(X_t^\mu,Y) \rangle_{H_{\Sigma_t}}
    \end{equation}
in distribution.
\end{corollary}
\begin{proof}
    Combining Lemmas \ref{lem: relation between CM spaces of U and Vt} $(iv)$ and \ref{lem: training loss}, it follows that
\begin{equation*}
    \langle x_0 , m_t(X_t^\mu, Y) \rangle_{H_{\Sigma_t}}
    =  \langle x_0, m_t(X_t^\mu, Y) \rangle_{H_{\widetilde{\C}_t}} + \langle x_0, A^*\Gamma^{-1}A m_t(X_t^\mu, Y) \rangle_H
\end{equation*}
in distribution. Therefore, we obtain
\begin{align*} 
    \langle x_0 &, m_t(X_t^\mu, Y) \rangle_{H_{\Sigma_t}} \\
    &=  \langle x_0, m_t(X_t^\mu, Y) \rangle_{H_{\widetilde{\C}_t}} + \langle x_0, A^*\Gamma^{-1}A m_t(X_t^\mu, Y) \rangle_H  \\
    & = {e^{t/2}} \langle x_0, X_t^\mu \rangle_{H_{\widetilde{\C}_t}} +  \langle x_0 , \widetilde{\C}_t A^*C_t^{-1}(Y-e^{t/2}AX_t^\mu) \rangle_{H_{\widetilde{\C}_t}} + e^{t/2} \langle x_0, A^*\Gamma^{-1}AX_t^\mu \rangle_{H} \\
    & \quad\quad + \langle x_0, (e^t-1) A^*\Gamma^{-1}A \C A^* C_t^{-1} (Y-e^{t/2}AX_t^\mu) \rangle_H\\
    &= e^{t/2} \langle x_0, X_t^\mu \rangle_{H_{\widetilde{\C}_t}} +  e^{t/2} \langle x_0, A^*\Gamma^{-1}AX_t^\mu \rangle_{H} + \langle x_0, A^*C_t^{-1}(Y-e^{t/2}AX_t^\mu) \rangle_H \\
    & \quad\quad + (e^t-1) \langle x_0, A^*\Gamma^{-1}A \C A^* C_t^{-1} (Y-e^{t/2}AX_t^\mu) \rangle_H\\
    &= e^{t/2} \langle x_0, X_t^\mu \rangle_{H_{\widetilde{\C}_t}} +  e^{t/2} \langle x_0, A^*\Gamma^{-1}AX_t^\mu \rangle_{H} 
    + \langle x_0, A^* \Gamma^{-1} (Y-e^{t/2}X_t^\mu)\rangle_H \\
    & = e^{t/2} \langle x_0, X_t^\mu \rangle_{H_{\widetilde{\C}_t}} + \langle x_0,  A^*\Gamma^{-1 }Y  \rangle_H
\end{align*}
in distribution, where we have used that
\begin{eqnarray*}
    A^*C_t^{-1} + (e^t-1) A^*\Gamma^{-1}A \C A^* C_t^{-1} = A^*  \Gamma^{-1}\underbrace{\left(\Gamma + (e^t-1) A\C A^* \right)}_{=C_t}C_t^{-1} = A^*\Gamma^{-1}.
\end{eqnarray*}
This completes the proof.
\end{proof}


\begin{proof}[Proof of Proposition \ref{prop:identity_of_expectations}]
    Let $x_0 \in H_{\widetilde{\C}_t} \subset H$ where $\widetilde{\C}_t = (e^t-1)\C$. By virtue of Lemma \ref{lem: tranisition kernels}, we will prove that
    \begin{eqnarray}\label{eq:transition-ident}
        \frac{\int_H x_0 n_t(x_0, x) \mu^y(dx_0)}{\int_H n_t(x_0, x) \mu^y(dx_0)} = \frac{\int_H x_0 \tilde n_t(x_0, m_t(x,y)) \mu(dx_0)}{\int_H \tilde n_t(x_0, m_t(x,y)) \mu(dx_0)},
    \end{eqnarray}
    for $(x,y) \in H \in \R^m$ a.e. in $\L(X_t^\mu, Y)$. 
    Notice that it suffices to show that %for $x_0 \in H_{\widetilde{\C}_t}$,
    \begin{eqnarray} \label{eq:Bayes-identity}
        n_t(x_0, x) \frac{d \mu^y}{d \mu}(x_0)= \frac{1}{\widetilde{Z}(y)}\tilde{n}_t(x_0, m_t(x,y)),
    \end{eqnarray}
    for $x_0 \in H$ $\mu$-a.e. and for $(x,y) \in H \in \R^m$ a.e. in $\L(X_t^\mu, Y)$, where $\widetilde{Z}(y)$ does not depend on $x_0$, hence will be canceled out in \Eqref{eq:transition-ident}. 
    By Bayes' theorem \Eqref{eq:posterior},
    \begin{eqnarray*}
        \frac{d \mu^y}{d \mu}(\cdot) =  \frac{1}{Z(y)}\exp\left( -\frac{1}{2}\norm{A\cdot-y}^2_\Gamma \right) \quad \text{in}\; L^1(\mu),
    \end{eqnarray*}
    Let us now write
    \begin{eqnarray*}
         -\frac{1}{2} \norm{Ax_0-y}^2_\Gamma 
        & = &  \langle x_0, A^*\Gamma^{-1}y \rangle_H - \frac{1}{2}\langle x_0, A^*\Gamma^{-1}Ax_0 \rangle_H  -\frac{1}{2} \norm{y}_\Gamma^2
    \end{eqnarray*}
   and set $\widetilde{Z}(y) = Z(y)\exp{\left(\frac{1}{2} \norm{y}_\Gamma^2\right)}$. For $x_0 \in H_{\widetilde{\C}_t}$,
    \begin{equation}
    \begin{aligned}
    \label{eq:part-of-prop-proof}
    & n_t(x_0, X_t^\mu)  \exp\left(\langle x_0, A^*\Gamma^{-1}Y \rangle_H - \frac{1}{2}\langle x_0, A^*\Gamma^{-1}Ax_0 \rangle_H  -\frac{1}{2} \norm{Y}_\Gamma^2\right) \\
    &= \frac{1}{\widetilde{Z}(Y)} \exp\big( e^{t/2} \langle x_0, X_t^\mu \rangle_{H_{\widetilde{\C}_t}} - \frac{1}{2}\norm{x_0}_{H_{\widetilde{\C}_t}}^2 + \langle x_0, A^* \Gamma^{-1}Y \rangle_H - \frac{1}{2} \langle x_0, A^* \Gamma^{-1} A x_0 \rangle_H \big) \\
    &=  \frac{1}{\widetilde{Z}(Y)}\exp\left( e^{t/2} \langle x_0, X_t^\mu \rangle_{H_{\widetilde{\C}_t}} + \langle x_0, A^* \Gamma^{-1} Y \rangle_H - \frac{1}{2}\norm{x_0}_{H_{\widetilde{\C}_t}}^2 -\frac{1}{2} \langle x_0, A^* \Gamma^{-1} A x_0 \rangle_H \right) \\ 
    &= \frac{1}{\widetilde{Z}(Y)}\exp\left( \langle x_0, m_t(X_t^\mu,Y) \rangle_{H_{\Sigma_t}} - \frac{1}{2} \norm{x_0}_{H_{\Sigma_t}}^2\right),
    \end{aligned}
    \end{equation}
    in distribution, where we have used Corollary \ref{cor:ID1} and the identity
    \begin{equation}\label{eq:ID2}
        \norm{x_0}_{H_{\widetilde{\C}_t}}^2 + \langle x_0, A^*\Gamma^{-1}Ax_0 \rangle_{H} = \norm{x_0}_{H_{\Sigma_t}}^2,
    \end{equation}
    which follows from Lemma \ref{lem: relation between CM spaces of U and Vt} $(iv)$.
    
    Now by Lemma \ref{lem: relation between CM spaces of U and Vt} $(ii)$, $x_0\in H_{\Sigma_t}$ and thus the Cameron--Martin theorem gives
    \begin{eqnarray*}
        \exp\left( \langle x_0, m_t(X_t^\mu,Y) \rangle_{H_{\Sigma_t}} - \frac{1}{2}\norm{x_0}^2_{H_{\Sigma_t}} \right) = \tilde{n}_t(x_0, m_t(X_t^\mu,Y))
    \end{eqnarray*}
    in distribution,     which together with \Eqref{eq:part-of-prop-proof} lead to \Eqref{eq:Bayes-identity}. Therefore, \Eqref{eq:transition-ident} holds true, which completes the proof.
\end{proof}

\subsection{Gaussian example}
\begin{proof}[Proof for Lemma \ref{lem: score function is compact}]
     Let $\widetilde \C$ be a covariance operator such that $\mu\left(\widetilde{\C}^{1/2}(H)\right) = 1$. As $S_0^{1/2}(H)$ is the intersection of all linear subspaces of full measure under $\mu$ (see Prop. 4.45 in \citealp{hairer2023introductionSPDE}), it holds $S_0^{1/2}(H) \subset \widetilde{\C}^{1/2}(H)$. 

     We now find another covariance operator $\C$ such that the score function $\tilde s$ corresponding to $\C$ is bounded linear.
     Let $\C$ be such that $\C^{1/2}(X) \subset S_0^{1/2}(X)$ for any linear subset  $X$ of $H$.
    This implies
    \begin{eqnarray}\label{eq: inclusion of range of covariances}
        S_0(H) \subset \C(H) = \Sigma_t(H).
    \end{eqnarray}
    %Hence the map $S_0^{-1} \Sigma_t$ is well defined and bounded ($S_0: H \rightarrow S_0(H) \subset \Sigma_t(H)$ is invertible).
   We identify the law of $\widetilde{X}_t^{\mu}$ by using the relation \begin{eqnarray*}
        \widetilde{X}_t^{\mu}|x_0 \sim \N(x_0, \Sigma_t),
    \end{eqnarray*}
    where $x_0$ is a realisation of the prior $\mu = \N(0, S_0)$. Hence
    \begin{eqnarray*}
        \widetilde{X}_t^{\mu} \sim \N(0, \Sigma_t + S_0).
    \end{eqnarray*}
    By the reasoning of Lemma 4.4 in \citet{hairer2006analysis} it holds that 
    \begin{eqnarray*}
        \widetilde{X}_0^{\mu} |\widetilde{X}_t^{\mu} \sim \N(m', C')  
    \end{eqnarray*}
    with some covariance operator $C'$ and
    \begin{eqnarray*}
        m' &=& \Sigma_t (\Sigma_t + S_0)^{-1} \widetilde{X}_t^{\mu}. 
    \end{eqnarray*}
    Note that in Lemma 4.4 of \citet{hairer2006analysis}, the operator $\Sigma_t (\Sigma_t + S_0)^{-1}$ is defined as a measurable extension of the bounded map
    \begin{eqnarray*}
        A: (\Sigma_t + S_0)^{1/2}(H)\rightarrow H, \qquad x \mapsto \Sigma_t (\Sigma_t + S_0)^{-1}x
    \end{eqnarray*}
    to the whole space $H$ as per Theorem II.3.3 in \citet{DaleckijJurijL1991Made}. In our case, it is possible to give an explicit formula for a possible extension using the inclusion \Eqref{eq: inclusion of range of covariances}. %, which is not assumed in the context of \cite{hairer2006analysis}.
    We define
    \begin{eqnarray*}
        \Sigma_t (\Sigma_t + S_0)^{-1} &:=& \left( (\Sigma_t + S_0)^{-1} \Sigma_t \right)^* \\
        & = & \left( (I+ \Sigma_t^{-1}S_0)^{-1} \Sigma_t^{-1}S_0 \right)^*.
    \end{eqnarray*}
   This map coincides with $A$ on $(\Sigma_t + S_0)^{1/2}(H)$ and as we will now show, it is defined even on the whole space $H$.
   The operator $\Sigma_t: H \rightarrow \Sigma_t(H)$ is  bounded and invertible. Hence  also $\Sigma_t^{-1}: \Sigma_t(H) \rightarrow H$ is  bounded. By \Eqref{eq: inclusion of range of covariances} the map $\Sigma_t^{-1} S_0$ is well defined and bounded. 
    % Using this definition it is possible to repeat the proof of \cite{hairer2006analysis} to avoid the implicit extension. (Another line of thought: we have found a suitable extension) Compare with exercise 3.72 in \cite{hairer2023introductionSPDE} for a slightly weaker version of the extension.
    We can now identify the score function $\tilde s$ by using the previous equality for the conditional expectation
    \begin{eqnarray*}
        \tilde s(z, t; \mu) &=& -\left(z- \E(\widetilde{X}^{\mu}_0 |\widetilde{X}^{\mu}_t = z)\right) \\
        & = &-z + S_0 (\Sigma_t + S_0)^{-1}z \\
        & = & \left[S_0 (\Sigma_t + S_0)^{-1} - I\right]z. \\
        & = &- \Sigma_t (\Sigma_t + S_0)^{-1}z.
        %& =& -\Sigma_t S_0^{-1}(\Sigma_t S_0^{-1} + I )^{-1}z.
    \end{eqnarray*}
    This yields the claim.
    \end{proof}

\subsection{Conditional score matching}\label{sec: proof of conditional score matching}

Recall that $m_t(x, y) = e^{t/2}x + (e^t - 1)\C A^* C_t^{-1}(y - e^{t/2}Ax)$ and $\xi_t(x,y) :=  \C A^* \Gamma^{-1} y + \lambda(t) x$ for $x\in H$ and $y\in \R^m$. Moreover, $R_t := (e^t-1) I + (e^t-1)^2 \C A^* C_t^{-1} A : H\to H$.

\begin{lemma}\label{lem: writing m_t for no forward}
The operator $R_t : H\to H$ is bijective and the identity
\begin{eqnarray*}
    m_t(x,y) = R_t \xi_t(x, y).
\end{eqnarray*}
holds for all $x\in H$ and $y\in \R^m$.
\end{lemma}
\begin{proof}
    First note that bijectivity is implied by Lemma \ref{lem: Useful equality that show that Sigma is still p.d.} (b).
    By direct computation,
    \begin{eqnarray*}
        R_t \xi_t(x, y)&= & R_t \left(\C A^* \Gamma^{-1}y + \lambda(t) x\right) \\ 
        & = & (e^t-1) \C A^* \Gamma^{-1}y - (e^t-1)^2  \C A^* C_t^{-1} A \C A^* \Gamma^{-1}y \\
        & & + e^{t/2} x - e^{t/2}(e^t-1)\C A^* C_t^{-1}Ax \\
        & = & e^{t/2}x - e^{t/2}(e^t-1)\C A^* C_t^{-1}Ax \\
        & & + (e^t-1) \C A^* (C_t^{-1} [ \underbrace{(e^t-1) A \C A^*+ \Gamma}_{= C_t} ] + (e^t-1)C_t^{-1} A \C A^* ) \Gamma^{-1}y \\
        & = & e^{t/2}x + (e^t-1) \C A^* C_t^{-1} \left( y- e^{t/2}Ax\right) \\
        & = & m_t(x, y).
    \end{eqnarray*}
    This completes the proof.
\end{proof}

\begin{proof}[Proof of Lemma \ref{lem: conditional score matching}]
We replicate the arguments of \cite{baldassari2024conditional} adapted to our setting. First, observe that 
\begin{multline*}
     \norm{r_\theta(R_t^{-1} \tilde x_t, t; \mu) - \tilde s(\tilde x_t, t; \mu) - \tilde x_t}_H^2  \\=  \norm{r_\theta(R_t^{-1} \tilde x_t, t; \mu)}_H^2 
    +  \norm{ \tilde s(\tilde x_t, t; \mu) - \tilde x_t}_H^2 - 2 \langle r_\theta(R_t^{-1} \tilde x_t, t; \mu) , \tilde s(\tilde x_t, t; \mu) - \tilde x_t \rangle_H.
\end{multline*}
    It holds by definition \Eqref{eq: posterior score in infinite dimensional space}
    \begin{align*}
        \E_{\tilde x_t \sim \L(\widetilde{X}_t^{\mu})} \langle & r_\theta(R_t^{-1}\tilde x_t, t; \mu) , \tilde s(\tilde x_t, t; \mu) - \tilde x_t \rangle_H  \\
        &= - \E_{\tilde x_t \sim \L(\widetilde{X}_t^{\mu})} \left\langle r_\theta(R_t^{-1}\tilde x_t, t; \mu) ,\E_{\tilde x_0 \sim \L(\widetilde{X}_0^{\mu} | \tilde x_t)} (\tilde x_t - \tilde x_0) - \tilde x_t \right\rangle_H \\
        & = \E_{\tilde x_0 \sim \L(\widetilde{X}_0^{\mu})} \E_{\tilde x_t \sim \L(\widetilde{X}_t^{\mu} | \tilde x_0 )} \langle r_\theta(R_t^{-1}\tilde x_t, t; \mu), \tilde x_0 \rangle_H.
    \end{align*}
    Hence it holds
    \begin{multline*}
        \lambda(t)^2 \E_{\tilde x_t \sim \L(\widetilde{X}_t^{\mu})} \norm{r_\theta(R_t^{-1} \tilde x_t, t; \mu) - \tilde s(\tilde x_t, t; \mu) + \tilde x_t}_H^2 \\
        =  V'(t) + \lambda(t)^2\E_{\tilde x_0 \sim \L(\widetilde{X}_0^{\mu})} \E_{\tilde x_t \sim \L(\widetilde{X}_t^{\mu} | \tilde x_0 )} \norm{r_\theta(R_t^{-1} \tilde x_t, t; \mu)- \tilde x_0 }_H^2
    \end{multline*}
    with
    \begin{equation}
        \label{eq:Vdash_const}
        V'(t) = \lambda(t)^2 \E_{\tilde x_t \sim \L(\widetilde{X}_t^{\mu})} \norm{\tilde s(\tilde x_t, t, \mu) - \tilde x_t}_H^2 - \lambda(t)^2 \E_{\tilde x_0 \sim \L(\widetilde{X}_0^{\mu})} \norm{\tilde x_0}_H^2.
    \end{equation}
    To conclude, we add expectation with respect to $t \sim [\delta, T]$.
    Note that by Assumption \ref{ass: stability estimate; no forward process} the first term on the rhs of \Eqref{eq:Vdash_const} is uniformly bounded in $t$ and by elementary calculations
    \begin{eqnarray*}
        \E_{\tilde x_0 \sim \L(\widetilde{X}_0^{\mu})} \lambda(t)^2 \norm{ \tilde x_0}_H^2 
         \le  \lambda(\delta)^2 \E_{x\sim \mu}\norm{x}_H,
    \end{eqnarray*}
    such that
    \begin{eqnarray*}
        V:=\E_{t \sim [\delta, T]} V'(t) \le \frac{1}{T-\delta}\sup_{t \in [\delta, T]} V'(t) < \infty.
    \end{eqnarray*}
    This concludes the proof.

\end{proof}
    
\section{Proofs for Section \ref{sec: Estimating numerical errors}}


\begin{proof}[Proof of Theorem \ref{thm: stability estimate}]
Below, we use the notation $f\lesssim g$, if $f(x) \leq C g(x)$ for all $x$ with some universal constant $C>0$. 
Recall that true solution $W_t^{\mu^y}$ of the time-reversed denoising process and corresponding approximative solution process $V_t^{\mu^y}$ satisfy
\begin{align*}
    dW_t^{\mu^y} &= \left(-\frac{1}{2} W_t^{\mu^y} - s(W_t^{\mu^y}, T-t; \mu^y) \right) dt + \C^{1/2} dB_t, \\
    d V_t^{\mu^y} &= -\frac{1}{2}  V_\low{t}^{\mu^y} - \lambda(T-\low{t})\left[ r_\theta\left(\xi_{T - \low t}(V_{\low{t}}^{\mu^y}, y), T-\low{t}; \mu\right)  - e^{(T-\low{t})/2} V_{\low{t}}^{\mu^y} \right] dt +\C^{1/2} dB_t.
\end{align*} 
In what follows, we abbreviate the expectation 
$\E_{y \sim \pi_y}\E_{(w_T, v_{T-\delta}) \sim \L( W_T^{\mu^y},V_{T-\delta}^{\mu^y})}$ as $\E$, unless otherwise specified.

{\bf Decomposition of the error.} Let us first consider the difference
\begin{eqnarray*}
    \E \norm{w_t -  v_t}_H^2 & = & \E \norm{w_t - w_0 - (v_t - v_0) + (w_0 - v_0)}_H^2\\
    &\lesssim&  \E \norm{w_t - w_0 - (v_t - v_0)}_H^2 + \varepsilon_{INIT},
\end{eqnarray*}
where $\varepsilon_{INIT} := \E \norm{w_0 - v_0}_H^2$. 
Now we decompose the difference into three terms as follows:
\begin{equation*}
	w_t - w_0 - (v_t - v_0) = \int_0^t \left[{\mathcal I}_1(\tau) + {\mathcal I}_2(\tau) + {\mathcal I}_3(\tau)\right]d\tau,
\end{equation*}
where we have
\begin{eqnarray*}
	{\mathcal I}_1(\tau) & = & -\frac 12 \left(w_\tau - w_{\low \tau}\right) +
	s(w_\tau, T-\tau; \mu^y)- s(w_{\low\tau}, T-\low \tau; \mu^y), \\
	{\mathcal I}_2(\tau) & = & s(w_{\low \tau}, T - \low \tau; \mu^y) - \lambda(T-\low \tau)\left(r_\theta(\xi_{T-\low\tau}(w_{\low\tau}, y), T-\low \tau; \mu) -e^{(T-\low \tau)/2}w_{\low\tau}\right) \quad \text{and}\\
	{\mathcal I}_3(\tau) & = & \lambda(T-\low \tau) \left\{r_\theta(\xi_{T-\low\tau}(w_{\low\tau}, y), T - \low \tau; \mu) - r_\theta(\xi_{T-\low\tau}(v_{\low\tau}, y), T - \low \tau; \mu)\right\} \\
	& & - (\lambda(T-\low \tau) e^{(T-\low \tau)/2}+1)\left(w_{\low\tau} - v_{\low\tau} \right).
\end{eqnarray*}

{\bf Bound for $\varepsilon_{INIT}$.} Recall that $v_0 =  (1-e^{-T}) Z$ and $w_0 = (1-e^{-T}) Z + e^{-T/2} X_0$, where $Z \sim \N(0,\C)$ and $X_0\sim \mu^y$. It directly follows that
\begin{eqnarray*}
    \varepsilon_{INIT} \leq  e^{-T} \E_{y\sim\pi_y}\E_{x\sim \mu^y} \norm{x}_H^2 = e^{-T} \E_{x\sim \mu} \norm{x}_H^2.
\end{eqnarray*}
where we applied marginalization of the joint distribution.

{\bf Contribution from ${\mathcal I_1}$}: We observe that
\begin{align*}
    & \varepsilon_{NUM} := \E \int_0^{T-\delta} \norm{{\mathcal I}_1(\tau)}_H^2 d\tau \\
    & = \;\E \int_0^{T-\delta} \norm{-\frac 12 \left(w_\tau - w_{\low\tau}\right) +
	s(w_\tau, T-\tau; \mu^y)- s(w_{\low\tau}, T-\low \tau; \mu^y)}_H^2 d \tau \\
    & \lesssim \E \int_0^{T-\delta} \norm{w_\tau - w_{\low\tau}}_H^2 d\tau + \E \int_0^{T-\delta} \norm{s(w_\tau, T-\tau; \mu^y)- s(w_{\low\tau}, T-\low \tau; \mu^y)}_H^2 d\tau.
\end{align*}
By Lemma 2 in \citet{pidstrigach2023infinitedimensional}, it holds that
\begin{eqnarray*}
    s(X_t^{\mu^y}, t; \mu^y) = e^{(t- \tau)/2}\E\left( s(X_\tau^{\mu^y}, \tau; \mu^y) | X_t^{\mu^y}\right), \quad 0 < t \le \tau \le T,
\end{eqnarray*}
for $\pi_y$-a.e. $y\in\R^m$. Therefore, we deduce by Lemma 11 in \citet{chen2023improved} that for the time-reversed process it holds that
\begin{align*}
     \E \big\|& s(w_\tau, T-\tau; \mu^y)   - s(w_{\low\tau}, T-\low{\tau}; \mu^y)\big\|_H^2  \\
     & \le 4 \E \norm{s(w_\tau, T-\tau; \mu^y) - s( e^{(\tau - \low{\tau})/2}w_{\low\tau}, T- \tau; \mu^y)}_H^2 \\
     & \quad\quad+ 2 \left(1-e^{\tau - \low{\tau}}\right)^2 \E \norm{s(w_\tau, T-\tau; \mu^y)}_H^2 \\
     & \leq 4 L_s^2(T-\tau) \E \norm{w_\tau- e^{(t - \low{t}) / 2}w_{\low{\tau}}}_H^2 + 2(1-e^{\tau - \low{\tau}})^2 \E \norm{s(w_\tau, T - \tau; \mu^y)}_H^2,
\end{align*}
where we used the Lipschitz continuity of $s$.
We note that 
\begin{eqnarray}\label{eq: transition kernel ou process}
     w_s | w_t \sim \N\left(e^{-(t - s)/2} w_t, (1-e^{-(t - s)} ) \C\right),
\end{eqnarray}
for $T \ge t \ge s \ge 0$, see \citet{FactsForInfiniteDimOU}.
It immediately follows that 
\begin{eqnarray*}
    \E \norm{w_\tau - e^{(\tau - \low{\tau}) / 2}w_{\low{\tau}}}_H^2 = (1-e^{-(\tau - \low{\tau})}) \tr_H(\C),
\end{eqnarray*}
and, consequently,
\begin{align*}
    \E \big\|w_\tau - & w_{\low\tau}\big\|_H^2 \lesssim \E \norm{w_\tau - e^{(\tau - \low{\tau}) / 2}w_{\low\tau}}_H^2 + (1-e^{(\tau - \low{\tau}) / 2})^2 \E \norm{w_{\low\tau}}_H^2 \\
    & \lesssim  (1-e^{-(\tau - \low{\tau})}) \tr_H(\C) \\
    & \quad\quad + (1-e^{(\tau - \low{\tau}) / 2})^2 \left(\E_{y\sim\pi_y}\E_{x\sim \mu^y} \norm{x}_H^2 + \left((1-e^{-(T-\low{\tau})}\right) \tr_H(\C)\right) \\
    & \lesssim (1-e^{-\Delta t}) \tr_H(\C)
    +(1-e^{\Delta t/ 2})^2 \left(\E_{x\sim\mu} \norm{x}^2_H + \tr_H (\C)\right).
\end{align*}
Combining the arguments yields
\begin{eqnarray*}
    \varepsilon_{NUM} & \lesssim &
    (T-\delta) \left((1-e^{-\Delta t}) \tr_H(\C)
    +(1-e^{\Delta t/ 2})^2 \left(\E_{x\sim\mu} \norm{x}^2_H + \tr_H (\C)\right)\right) \\
    & & + L_s^2(T-\tau) (1-e^{-\Delta t}) \tr_H(\C) 
    + (1-e^{\Delta t})^2 \E \norm{s(w_\tau, T - \tau; \mu^y)}_H^2.
\end{eqnarray*}
Since the last expectation and $\E_{x\sim\mu} \norm{x}_H^2$ are bounded by assumption, and note that $1-e^{-\Delta t} = {\mathcal O}(\Delta t),\; (1-e^{\Delta t/ 2})^2 = {\mathcal O}(\Delta t)$, we obtain the required upper bound.


{\bf Contribution from ${\mathcal I_2}$}: Applying Theorem \ref{thm:score_id1 infinite dimension}, we obtain
\begin{eqnarray*}
	{\mathcal I_2}(\tau) & = & \lambda(T-\low \tau)\big\{\Tilde{s}(R_{T-\low \tau} \xi(w_{\low\tau},y), T-\low \tau; \mu) + R_{T-\low \tau} \xi(w_{\low\tau},y)  \\
	& & \quad - r_\theta(\xi_{T-\low\tau}(w_{\low\tau}, y), T - \low \tau; \mu) \big\} \\
	& = & \lambda(T-\low \tau)\big\{\Tilde{s}(m_{T-\low \tau} (w_{\low\tau},y), T-\low \tau; \mu)  + m_{T-\low \tau} (w_{\low\tau},y) \\
	& &   - r_\theta(R_{T-\low \tau}^{-1} m_{T-\low \tau}(w_{\low\tau}, y), T - \low \tau; \mu) \big\}
\end{eqnarray*}
and now it follows by reversing time and applying Lemma \ref{lem: training loss} that
\begin{equation*}
	\E \int_0^{T-\delta} \norm{{\mathcal I}_2(\tau)}_H^2 d\tau = \varepsilon_{LOSS},
\end{equation*}
where $\varepsilon_{LOSS}$ is given by \Eqref{eq:discrete_loss}.

{\bf Contribution from ${\mathcal I_3}$}: By triangle inequality and the assumption on uniform Lipschitzness of $r_\theta$ we have
\begin{align*}
	\big\|{\mathcal I}_3 & (\tau)\big\|_H \leq  
	\lambda(T-\low\tau) L_s(T-\low\tau)
	\norm{\xi_{T-\low\tau}(w_{\low\tau}, y) -\xi_{T-\low\tau}(v_{\low\tau}, y)}_H \\
%	+(\lambda(T-\low\tau)e^{(T-\low\tau)/2} +1)\norm{W_{\low \tau}-V_{\low \tau}}_H\bigg) \\
	&   %\lambda(T-\low\tau)\bigg(L_s(T-\low\tau)\lambda(T-\low\tau)
	+\lambda(T-\low\tau)e^{(T-\low\tau)/2} +1) \norm{w_{\low\tau}-v_{\low\tau}}_H \\
	& \le \kappa_s(T-\low\tau) \norm{w_{\low\tau}-v_{\low\tau}}_H,
\end{align*}
where we abbreviate $\kappa_s(\tau') = L_s(\tau')\lambda(\tau')^2 +\lambda(\tau')e^{\tau'/2} +1$ for convenience. Note that by assumption $\kappa_s(\cdot) \in L^2[\delta, T]$.


Combining the estimates, we obtain
\begin{equation*}
	\E \norm{w_{T-\delta} -  v_{T-\delta}}_H^2 \lesssim
	\varepsilon_{NUM} + \varepsilon_{LOSS} + \varepsilon_{INIT} + \E \int_0^{T-\delta} \kappa_s(T-\low\tau)^2 \norm{w_{\low\tau}-v_{\low\tau}}_H^2 d\tau.
\end{equation*}
Applying Gr\"onwall's inequality, it follows
\begin{equation*}
	\E \norm{w_{T-\delta} -  v_{T-\delta}}_H^2 \lesssim
	\left(\varepsilon_{NUM} + \varepsilon_{LOSS} + \varepsilon_{INIT}\right) \exp\left(\int_0^{T-\delta} \kappa_s(T-\low\tau)^2 d\tau\right).
\end{equation*}


We may factor in the truncation by utilizing \Eqref{eq: transition kernel ou process}
\begin{eqnarray*}
    \E {\norm{w_T -  w_{T-\delta}}_H^2} &=& \E_{y\sim \pi_y}\E_{x\sim\mu^y, z\sim \N(0,\C)}\norm{(1- e^{-\delta/2})x + \sqrt{1-e^{-\delta}} z}_H^2 \\
    & \lesssim &  (1- e^{-\delta/2})^2 \E_{x\sim\mu}\norm{x}_H^2 + (1-e^{-\delta})  \tr_H(\C) \\
    & = & {\mathcal O}(\delta).
\end{eqnarray*}
Hence, we conclude
\begin{equation*}
    \E \norm{w_T -  v_{T-\delta}}_H^2
    \lesssim \left(\varepsilon_{NUM} + \varepsilon_{LOSS} + \varepsilon_{INIT}\right) \exp\left( 2 \int_0^{T-\delta} \kappa_s(T-\low\tau)^2 d\tau\right) +  \delta,
\end{equation*}
which yields the result.
\end{proof}





\section{Details of numerical implementation}\label{section: Details of NN}

\paragraph{Neural network architecture}
In Figure \ref{fig: FNO}, we highlight the FNO architecture that we use to approximate the score function in our approach. Here, we implement 4 hidden layers and $s$ denotes the number of pixels in both horizontal and vertical directions, while $h$ represents the number of hidden nodes. The Conditional method employs the same architecture but omits the transform $f(t)$ and includes an additional input node $y$. This modification affects only the dimensionality of the first layer, changing it to $\R^{s \times s \times 3}$, while all subsequent layers remain unchanged. Consequently, the total number of parameters remains in the same order. 
The unconditional score approximation also follows the same architecture as UCoS but without the transform $f(t)$.

\begin{figure}[h]
    \centering
    \includegraphics[width=0.8\linewidth]{Images/FNO/FNO.pdf}
     \caption{Our FNO architecture. The scalar multiplication is given by $f(t)x = \frac{x}{1+ \text{std}(x)}$.} Activation function $\sigma$ is given by Relu. The lifting block $\mathcal{R}$ is given by a linear layer and $\mathcal{Q}$ by two linear layers and an activation function. The Fourier layers are followed by a batch--normalization layer before the activation function. Below each layer is the dimensionality of a tensor after passing through each layer
    \label{fig: FNO}
\end{figure}

We use the FNO architecture from \citet{baldassari2024conditional} and choose $\delta = 5 \cdot 10^{-3}$ and $T = 1$. Similar to the works of \citet{pidstrigach2023infinitedimensional, baldassari2024conditional}, we run the forward SDE with a non-constant speed function leading to the SDE
\begin{eqnarray*}
    d X_t = -\frac{1}{2} \beta(t) X_t \dd t + \sqrt{\beta(t) \C} \dd W_t
\end{eqnarray*}
with $\beta(t) = 0.05 + t (10-0.05)$.
%From this, we train using the ADAM optimizer with learning rate $10^{-4}$ and ReduceLROnPlateau scheduler to 

We train the neural network for 10 epochs for the 32-nodes architecture, 30 epochs for the 64-nodes architecture and 65 epoch for the 128-nodes architecture, where training is done on a Nvidia A100 GPU with 80 GB of memory. Training takes from 14-46 minutes for the unconditional method (depending on the architecture), 17-47 minutes for the conditional method and 24-51 minutes for UCoS.

During the training process we use the Adam stochastic optimizer with linearly decaying learning rate ranging from 0.002 to 0.0005. 
%While we monitor the loss function \Eqref{eq:discrete_loss}, gradients are computed without scaling by $\lambda(t)$.
Samples are generated on the same machine.
A uniform Euler--Maruyama approximation with 1000 steps is employed. 
Additional information can be found in the numerical implementation \ExplSyntaxOn
\bool_if:NTF \accepted
    {\url{https://github.com/FabianSBD/SBD-task-dependent/tree/UCoS}.}
    {\url{https://github.com/FabianSBD/SBD-task-dependent/tree/UCoS}.}
\ExplSyntaxOff




\paragraph{Comparison methods}
For clarity, let us give the precise score approximation of the following methods that we compare our method with. Note that some of these methods are used in combination with some other sampling method in the reference but we will always utilise backwards in time Euler--Maruyama approximations.
\begin{itemize}
%     \item SDE LD \cite{jalal2021robustannealedlangevian} adjusts the unconditional score function by adding the gradient of the negative log-likelihood: 
% \begin{eqnarray*}
%     s_\theta(x, t; \mu^y ) = s(x, t, \mu) + A^* \Gamma^{-1} (y-Ax).
% \end{eqnarray*}
% This approximation is exact only at $t = 0$.
\item SDE ALD uses the score approximation of \citet{jalal2021robustannealedlangevian}, which is given by
\begin{eqnarray}\label{eq: SDE ALD}
    s_\theta(x, t; \mu^y ) = s(x, t, \mu) + A^* (\Gamma + \gamma_tI)^{-1}(y-Ax)
\end{eqnarray}
for the hyper-parameter $\gamma_t$ under the assumption $\Gamma = \sigma^2 I $ for some $\sigma^2 > 0$.
In line with \citet{feng2023scorebased}, we tune $\gamma_t$ such that the additive term has equal norm to the score function. 
\item DPS \cite{chung2023diffusion} employs a similar idea to SDE ALD by removing the hyper--parameter $\gamma_t$ and changing the mean of the Gaussian likelihood to be an estimate of $x_0$:
\begin{eqnarray*}
    s_\theta(x, t; \mu^y ) = s(x, t, \mu) - \rho\nabla_x \norm{y- A(\hat x_0(x))}_2^2
\end{eqnarray*}
for some $\rho$ which is chosen such that $\rho =\xi / \norm{y- A(\hat x_0(x))}$ for some constant $\xi$. We use a grid-search algorithm to find the optimal $\xi$.
Above $\hat x_0(x)$ is an estimate of $\E(X_0 | X_t = x)$ using the definition of the score function \Eqref{def: definition of OU Score function in infinite dimension}.
\item Proj \cite{Dey_2024}: This projection-based approach adds a data consistency step before every reverse time Euler--Maruyama step:
\begin{eqnarray*}
    x_t &=& (\lambda A^\top A + (1-\lambda)I)^{-1} \left( (1-\lambda )x_t' + \lambda A^\top y_t\right). %\\
    %x_s' &=& \frac{1}{2} (s-t) x_t - s(x_t, t, \mu) + W_t - W_s
\end{eqnarray*}
We tune the hyper-parameter $\lambda$ by a grid search algorithm. Since the operator $(\lambda A^\top A + (1-\lambda)I)^{-1}$ is very large in terms of GPU-memory, we employ an iterative scheme with 10 iterations at each time step to solve the inverse (CT imaging problem) or a very small batch-size (Deblurring problem). This significantly increases runtime compared to applying a precomputed operator.

\item Conditional \cite{baldassari2024conditional}: This approach approximates the conditional score function
\begin{eqnarray*}
    s(x, t; \mu^y) = -\frac{1}{1-e^{-t}}\left( x- e^{t/2}\E (X_0^\mu | X_t^\mu = x, Y = y)\right)
\end{eqnarray*}
directly.
\end{itemize}


% \paragraph{Deblurring}
% The forward operator $A$ corresponds to a two-dimensional convolution with a Gaussian kernel over the image domain. We observe the convoluted function on the same grid with additive Gaussian noise. 

% with a Gaussian kernel with variance $\sigma^2 = 6$. The images are Observational noise is drawn with variance $5\%$ of the largest singular value of forward operator.

% \paragraph{CT imaging}
% As linear forward operator we choose sparse-view CT imaging with a maximum view angle of first 30 degrees. Observational noise is drawn with variance $5\%$ of the largest singular value of forward operator.

%%%%%%%%%%%%%%%%%%%%%%%%%%%%%%%%%%%%%%%%%%%%%%%%%%%%%%%%%%%%%%%%%%%%%%%%%%%%%%%
%%%%%%%%%%%%%%%%%%%%%%%%%%%%%%%%%%%%%%%%%%%%%%%%%%%%%%%%%%%%%%%%%%%%%%%%%%%%%%%

\section{Additional figures}
\label{sec: additional samples}
More samples and in higher resolution for all methods can be found below in Figure \ref{fig:samples inpainting} for the inpainting problem, in Figures \ref{fig:samples 32}, \ref{fig:samples 64} and \ref{fig:samples 128} for the CT imaging problem and Figure \ref{fig:deblurring samples 128} for the Deblurring problem.

\begin{figure}[H]
    \centering
    \includegraphics[width=\linewidth]{Images/Inpainting/samples3210.png}
    \caption{Posterior samples for Inpainting problem. Methods from top to bottom: SDE ALD, DPS, Proj, Conditional, UCoS, true posterior}
    \label{fig:samples inpainting}
\end{figure}


\begin{figure}[H]
    \centering
    \includegraphics[width=\linewidth]{Images/CT_imaging/32_nodes/samples3210.png}
    \caption{Posterior samples for CT imaging problem with a FNO architecture that uses 32 nodes per layer. Methods from top to bottom: SDE ALD, DPS, Proj, Conditional, UCoS.}
    \label{fig:samples 32}
\end{figure}

\begin{figure}[H]
    \centering
    \includegraphics[width=\linewidth]{Images/CT_imaging/64_nodes/samples6410.png}
    \caption{Posterior samples for CT imaging problem with a FNO architecture that uses 64 nodes per layer. Methods from top to bottom: SDE ALD, DPS, Proj, Conditional, UCoS.}
    \label{fig:samples 64}
\end{figure}

\begin{figure}[H]
    \centering
    \includegraphics[width=\linewidth]{Images/CT_imaging/128_nodes/samples12810.png}
    \caption{Posterior samples for CT imaging problem with a FNO architecture that uses 128 nodes per layer. Methods from top to bottom: SDE ALD, DPS, Proj, Conditional, UCoS.}
    \label{fig:samples 128}
\end{figure}

\begin{figure}[H]
    \centering
    \includegraphics[width=\linewidth]{Images/Deblurring/samples12810.png}
    \caption{Posterior samples for Deblurring problem. Methods from top to bottom: SDE ALD, DPS, Proj, Conditional, UCoS.}
    \label{fig:deblurring samples 128}
\end{figure}


Below in Figure \ref{fig:comparision-graph-big}, we reprint Figure \ref{fig: graph nodes vs error} in larger size.
\begin{figure}[H]
    \centering
    % First subfigure
    \begin{subfigure}[b]{0.49\linewidth} % Adjust width as needed
        \centering
        \includegraphics[width=\linewidth]{Images/CT_imaging/graph/graphbias.png}
    \end{subfigure}
    \hfill % Add spacing between figures
    % Second subfigure
    \begin{subfigure}[b]{0.49\linewidth} % Adjust width as needed
        \centering
        \includegraphics[width=\linewidth]{Images/CT_imaging/graph/graphstd.png}
    \end{subfigure}
    \caption{Comparison of the error dependence of the conditional method and UCoS on the parametrization of the score approximation. X-axis: number of nodes per layer in the FNO, Y-axis: L2 norm of the bias (left) or std (right), averaged over the number of pixels. Dashed lines correspond to $64 \times 64$ resolution and crosses correspond to the problem with $256 \times 256$ resolution.}
    \label{fig:comparision-graph-big}
\end{figure}



\end{document}
